%%%%%%%%%%%%%%%%%%%%%%%%%%%%%%%%%%%%%%%%%%%%%%%%%%%%%%%%%%%%%%%%%%%%%%%%%%%%
%% Author template for INFORMS Journal on Computing (ijoc)
%% Mirko Janc, Ph.D., INFORMS, mirko.janc@informs.org
%% ver. 0.95, December 2010
%%%%%%%%%%%%%%%%%%%%%%%%%%%%%%%%%%%%%%%%%%%%%%%%%%%%%%%%%%%%%%%%%%%%%%%%%%%%
%\documentclass[ijoc,blindrev]{informs3}
\documentclass[ijoc,blindrev]{informs3} % current default for manuscript submission
% If hyperref is used, dvi-to-ps driver of choice must be declared as
%   an additional option to the \documentclass. For example
%\documentclass[dvips,ijoc]{informs3}      % if dvips is used
%\documentclass[dvipsone,ijoc]{informs3}   % if dvipsone is used, etc.

\usepackage{tabularx}
\usepackage{adjustbox}

% --- Preamble split ---
\input{preamble/packages}
\input{preamble/macros}

% Graphics root
\graphicspath{{figures/}}

%%%%%%%%%%%%%%%%
\begin{document}
%%%%%%%%%%%%%%%%

% Line spacing / font size per INFORMS class
\OneAndAHalfSpacedXII  % current default line spacing
%%\OneAndAHalfSpacedXI
%%\DoubleSpacedXII
%%\DoubleSpacedXI

% Full title. Sample:
% \TITLE{Bundling Information Goods of Decreasing Value}
% Enter the full title:
\TITLE{BMSSPy: A Python Package and Empirical Comparison of Bounded Multi-Source Shortest Path Algorithm}

% Title or shortened title suitable for running heads. Sample:
% \RUNTITLE{Bundling Information Goods of Decreasing Value}
% Enter the (shortened) title:
\RUNTITLE{BMSSPy}

% Author's names for the running heads
% Sample depending on the number of authors;
% \RUNAUTHOR{Jones}
% \RUNAUTHOR{Jones and Wilson}
% \RUNAUTHOR{Jones, Miller, and Wilson}
% \RUNAUTHOR{Jones et al.} % for four or more authors
% Enter authors following the given pattern:
%\RUNAUTHOR{}
\RUNAUTHOR{Makowski et al.}

% Block of authors and their affiliations starts here:
% NOTE: Authors with same affiliation, if the order of authors allows, 
%   should be entered in ONE field, separated by a comma. 
%   \EMAIL field can be repeated if more than one author
\ARTICLEAUTHORS{%
\AUTHOR{Connor Makowski, Willem Guter, Tim Russell, Lucas Castro, Austin Saragih}
\AFF{Massachusetts Institute of Technology, Cambridge, MA, USA \EMAIL{\{conmak,wjguter,trussell,saragih\}@mit.edu}}
\AFF{Federal University of Amazonas, Manaus, AM, BRA \EMAIL{lucas.castro@icomp.ufam.edu.br}}

} % end of the block% end of the block

% --- Front matter ---
% Title, authors, history
% frontmatter/abstract.tex
\ABSTRACT{%
We introduce BMSSPy, {\color{blue}to our knowledge the first} open-source Python package that faithfully implements the Bounded Multi-Source Shortest Path (BMSSP) algorithm. This algorithm was recently discovered as a breakthrough that addresses Dijkstra’s sorting barrier for single-source shortest paths. BMSSPy implements the theoretical breakthrough with three practical {\color{blue}contributions}. First, {\color{blue}it preserves the graph transformations and data structures required for the algorithm's stated asymptotic guarantees while returning both distances and explicit paths}. Second, it clarifies pivot selection, tight-edge forest construction, and frontier management. Third, it handles graphs with non-unique path lengths, structural irregularities, unreachable nodes, and highly variable node degrees, {\color{blue}supporting robust execution on practical graph instances}.
{\color{blue}Experiments on synthetic gridgraphs and real geospatial networks show that BMSSP becomes increasingly competitive with general-purpose shortest-path implementations as graph size increases. On the largest tested gridgraphs, Python BMSSP is approximately $1.6$--$2.1\times$ faster than NetworkX Dijkstra and $2.4$--$4.2\times$ faster than igraph Dijkstra. A C++ implementation of BMSSP is roughly $6$--$8\times$ faster than its Python counterpart across most instances, highlighting the importance of implementation overhead. Nevertheless, both Python and C++ \texttt{SCGraph} Dijkstra remain faster throughout the observed benchmark range.}
{\color{blue}BMSSPy therefore provides a reproducible reference implementation and research platform for studying the practical behavior and further algorithm engineering of BMSSP and shortest-path algorithms.}
}%

% Sample 
%\KEYWORDS{deterministic inventory theory; infinite linear programming duality; 
%  existence of optimal policies; semi-Markov decision process; cyclic schedule}

% Fill in data. If unknown, outcomment the field
\KEYWORDS{Software, Algorithms, Shortest Path, Logistics, Distance Estimation}

% Example
%\HISTORY{%
%This paper was first submitted on January 1, 2025 and has been with the authors for X months for Y revision(s).
%}
%\HISTORY{}

% Abstract + keywords
% frontmatter/titleblock.tex

% Outcomment only when entries are known. Otherwise leave as is and 
%   default values will be used.
%\setcounter{page}{1}
%\VOLUME{00}%
%\NO{0}%
%\MONTH{Xxxxx}% (month or a similar seasonal id)
%\YEAR{0000}% e.g., 2005
%\FIRSTPAGE{000}%
%\LASTPAGE{000}%
%\SHORTYEAR{00}% shortened year (two-digit)
%\ISSUE{0000} %
%\LONGFIRSTPAGE{0001} %
%\DOI{10.1287/xxxx.0000.0000}%

% Title or shortened title suitable for running heads. Sample:
% \RUNTITLE{Bundling Information Goods of Decreasing Value}
% Enter the (shortened) title:




\maketitle

%%%%%%%%%%%%%%%%%%%%%%%%%%%%%%%%%%%%%%%%%%%%%%%%%%%%%%%%%%%%%%%%%%%%%%
% Samples of sectioning (and labeling) in IJOC
% NOTE: (1) \section and \subsection do NOT end with a period
%       (2) \subsubsection and lower need end punctuation
%       (3) capitalization is as shown (title style).
%
%\section{Introduction.}\label{intro} %%1.
%\subsection{Duality and the Classical EOQ Problem.}\label{class-EOQ} %% 1.1.
%\subsection{Outline.}\label{outline1} %% 1.2.
%\subsubsection{Cyclic Schedules for the General Deterministic SMDP.}
%  \label{cyclic-schedules} %% 1.2.1
%\section{Problem Description.}\label{problemdescription} %% 2.

% Text of your paper here
% Sections...
%%%%%%%%%%%%%%%%%%%%%%%%%%%%%%%%%%%%%%%%%%%%%%%%%%%%%%%%%%%%%%%%%%%%%%
\vspace{-20pt}
\section{Introduction}\label{sec:intro}
Solving the single-source shortest path (SSSP) problem on (un)directed graphs with non-negative weights is a foundational challenge in graph theory, operations research, theoretical computer science, and management science \citep{elmaghraby1970theory, orlin2010faster, duan_breaking_2025}. Given $m$ edges and $n$ nodes, the classic algorithm by \cite{dijkstra1959note}, enhanced by modern heap structures \citep{fredman1987fibonacci,chen2007priority}, achieves $O(m + n\log n)$ runtime. The term $O(n\log n)$ arises from the necessity of sorting distances—a limitation known as the \emph{sorting barrier}. 

This barrier is fundamental and makes Dijkstra’s algorithm uniquely optimal when putting nodes in the exact order of increasing distance \citep{haeupler2024universal}. However, recent breakthroughs reveal that this barrier is not essential if only distance values are required. The \emph{Bounded Multi-Source Shortest Path} (BMSSP) algorithm \citep{duan_breaking_2025} overcomes this, attaining $O(m\log^{2/3} n)$ complexity by eliminating global sorting and instead processing nodes within bounded distance ranges, resulting in an improved deterministic worst-case bound. This improvement is most consequential for sparse graphs, like many road networks.  %Given the arbitrary graphs and the need to convert to constant degree graphs, this is only true for sparse graphs


Although BMSSP marks a significant theoretical advance, {\color{blue}existing public implementations do not simultaneously provide a Python interface, the data-structural machinery required for the stated deterministic bound, and explicit treatment of several practical implementation details}. {\color{blue}This leaves an important gap in reproducibility and empirical evaluation.} Several details are left implicit in the theoretical descriptions, {\color{blue}including} constructing the tight-edge forest, handling ordering when path lengths are not unique, managing the suggested data structure, and coping with unreachable nodes or variable degree distributions. Therefore, a working implementation is essential to address these details, making the algorithm accessible and comparable to current practical methods, particularly in real-world transportation networks \citep{zhan1998shortest}. To that end, we have the following three research questions:
\begin{itemize}[topsep=0pt, partopsep=0pt, itemsep=0pt, parsep=0pt]
\item \textbf{RQ1:} How can BMSSP be implemented in a robust and reproducible manner?
\item \textbf{RQ2:} How can we resolve ambiguities that arise in practical implementations?
\item \textbf{RQ3:} How does BMSSP perform relative to both vanilla and {\color{blue}modern} Dijkstra variants on stylized grid graphs and real geospatial networks?
\end{itemize}

{\color{blue}
We address these questions with \texttt{BMSSPy}, to our knowledge the first fully open-source Python implementation of BMSSP that retains the algorithmic and data-structural components required for its stated deterministic complexity. After reviewing related literature and software in Section~\ref{sec:litproj}, the paper makes two main contributions:
}
\begin{itemize}

\item {\color{blue}
\textbf{Reproducible implementation and practical clarification.}
Sections~\ref{sec:pack} and~\ref{sec:func} develop an end-to-end implementation that returns shortest-path distances and explicit routes while making previously implicit components---including pivot selection, tight-edge forest construction, tie-breaking, set operations, and the underlying data structure---explicit and reproducible. This directly addresses RQ1 and RQ2.
}

\item {\color{blue}
\textbf{Empirical evaluation and scaling analysis.}
Section~\ref{sec:num} compares Python and C++ implementations of BMSSP with vanilla and {\color{blue}modern} Dijkstra variants on synthetic gridgraphs and real geospatial networks. We find that BMSSP becomes increasingly competitive with \texttt{NetworkX} and \texttt{igraph} as graph size grows, while \texttt{SCGraph} Dijkstra remains fastest. The benchmark and extrapolation analyses address RQ3, while the precision sensitivity analysis further supports the implementation choices examined under RQ2. Section~\ref{sec:conc} concludes and outlines directions for future work.
}

\end{itemize}

\section{Related Literature and Software}\label{sec:litproj}

We briefly review relevant literature in Section~\ref{subsec:rellit} and available software in Section~\ref{subsec:relproj}.

\subsection{Related Literature}\label{subsec:rellit}

The single-source shortest path (SSSP) problem on directed graphs with nonnegative weights is one of the most extensively studied problems in algorithms. The classical solution is Dijkstra’s algorithm \citep{dijkstra1959note}, which runs in \(O(m+n\log n)\) time when implemented with Fibonacci heaps \citep{fredman1987fibonacci}, relaxed heaps \citep{driscoll1988relaxed}, or modern priority queues \citep{chen2007priority}. Faster runtimes are possible when edge weights have additional structure. For undirected graphs with positive integer weights, \cite{thorup1999undirected} gives a linear-time \(O(m)\) algorithm. For directed graphs with positive integer weights, \cite{thorup2003integer} obtains \(O(m+n\log\log \min\{n,C\})\), where \(C\) is the maximum edge weight. If the graph contains only \(K\) distinct edge weights, \cite{orlin2010faster} devise a modified Dijkstra algorithm that runs in \(O(m)\) when \(nK\leq2m\) and in \(O\!\left(m\log \tfrac{nK}{m}\right)\) otherwise. In these approaches, the \(\log n\) factor is associated with ordering candidate distances, commonly referred to as the \emph{sorting barrier}.

Recent work clarifies when this barrier is unavoidable. \cite{haeupler2024universal} show that Dijkstra is \emph{universally optimal} when nodes must be output in increasing-distance order, even under powerful heap models. If this ordering requirement is removed, however, the barrier can be circumvented. Practical variants such as A* do not improve worst-case complexity, but can provide substantial speedups in structured settings such as road networks with informative heuristics \citep{hart1968formal}.

{\color{blue}
A substantial algorithm-engineering literature evaluates shortest-path methods empirically across synthetic and real networks. \citet{cherkassky1996shortest} compare multiple shortest-path algorithms across problem families, while the Ninth DIMACS Implementation Challenge established widely used benchmark formats, instances, and experimental protocols, including road networks with tens of millions of vertices and edges \citep{demetrescu2009shortest}. Road-network studies such as \citet{zhan1998shortest} identify Dijkstra-based variants among strong practical performers, and later surveys and implementations examine goal-directed and preprocessing methods, including Contraction Hierarchies and Transit Node Routing, on large transportation networks \citep{sommer2014shortest,bast2016route,makowski2025scgraph}. These studies motivate our use of complementary benchmark families: gridgraphs for controlled scaling and road, rail, and maritime networks for heterogeneous real-world evaluation. Our benchmark study is not intended to be exhaustive. This is because many advanced routing methods target point-to-point queries with substantial preprocessing. Instead, our direct comparison study focuses on Dijkstra-based baselines in software.
}


From the early survey of \cite{dreyfus1969appraisal} through the development of BMSSP, Dijkstra’s algorithm has remained a central general-purpose method for SSSP. When the output order of nodes is not required, however, \cite{duan_breaking_2025} show that the sorting barrier can be avoided. Their Bounded Multi-Source Shortest Path (BMSSP) algorithm attains deterministic \(O(m\log^{2/3} n)\) time by restricting relaxations to bounded distance ranges and reducing global sorting. BMSSP combines ideas from Bellman--Ford \citep{bellman1958routing,ford1956network} and Dijkstra. Earlier empirical work by \cite{golden1976shortest} also shows cases in which Bellman--Ford outperformed Dijkstra on sparse Euclidean graphs. In this context, BMSSP provides a deterministic SSSP approach that circumvents full sorting for directed and undirected graphs with nonnegative real weights. More recently, \citet{duan2026faster} further improve the BMSSP deterministic bound for directed real-weight SSSP to $O(m\sqrt{\log n}+\sqrt{mn\log n\log\log n})$. Our work focuses on the 2025 BMSSP construction because it is the algorithm implemented by \texttt{BMSSPy} and the basis of the implementation-engineering questions studied here.

{\color{blue}
Our work complements this literature by providing a reproducible BMSSP implementation and evaluating it against {\color{blue}various} Dijkstra implementations on synthetic and transportation networks.
}

\subsection{Related Software}\label{subsec:relproj}

{\color{blue}
Practical shortest-path performance also depends on implementation choices such as graph representation, programming language, priority-queue implementation, and software design. We therefore consider several widely used or {\color{blue}modern} implementations:
}

\begin{itemize}

\item \textbf{NetworkX:} A popular Python library for general graph analysis, offering Dijkstra's~\citep{dijkstra1959note} and A*~\citep{hart1968formal} implementations. It emphasizes usability and flexibility, although its pure-Python design can limit performance on large graphs~\citep{hagberg2008exploring}.

\item \textbf{igraph:} A high-performance graph library with C-backed algorithms and Python and R bindings~\citep{csardi2006igraph}. \texttt{igraph} provides efficient classical graph primitives but does not specifically target geospatial routing or implement BMSSP.

\item \textbf{SCGraph:} {\color{blue}A lightweight Python package for road, rail, and maritime shortest-path problems~\citep{makowski2025scgraph}. It extends real-network benchmarking to large global transportation networks and provides Dijkstra (including baseline Priority-heap by \cite{chen2007priority} and Buckets by \cite{dial1969buckets}) and A* implementations alongside preprocessing-based methods such as Contraction Hierarchies and Transit Node Routing.}

\end{itemize}

{\color{blue}
These packages provide practically relevant Python baselines but do not exhaust the broader high-performance shortest-path literature. Because our focus is general SSSP without preprocessing, we emphasize Dijkstra-based comparisons. Section~\ref{subsec:compar} also compares C++ implementations of BMSSP and \texttt{SCGraph} Dijkstra to measure the impact of Python.}

To our knowledge, no prior public implementation contains all components needed for BMSSP to realize its stated deterministic asymptotic guarantees. 
{\color{blue} The implementation reported in \citet{castro2025implementation} originally relied on hash-based structures and therefore provided expected rather than deterministic worst-case bounds. During the development of the present work, the implementation has continued to evolve, including incorporation of deterministic data-structure components. These developments illustrate the rapid evolution of BMSSP implementations and motivate our emphasis on reproducibility and explicit documentation of the components required for the deterministic bound.}
%\citet{castro2025implementation} provide a C++ prototype that relies on hash tables and therefore reports expected rather than worst-case running times. {\color{blue}This work demonstrates BMSSP's feasibility in a compiled language, but the hash-based structures do not retain the original deterministic worst-case guarantee. Our experiments therefore also include a C++ implementation equivalent to our Python BMSSP implementation, allowing us to separate language overhead from algorithmic overhead.}
\cite{valko_outperforming_2025} also recently released a Rust implementation evaluated on sparse payment-channel networks such as the Lightning Network. Their implementation omits several data-structural elements required for the theoretical bounds, resulting in a Dijkstra-equivalent worst-case runtime. %Their experiments report only small or negative speedups relative to Dijkstra and attribute this behavior to recursion overhead, constant factors, and memory-access costs.

{\color{blue}
Overall, these implementations highlight the distinction between reproducing BMSSP's core ideas and retaining the data structures required for its deterministic complexity bound. \texttt{BMSSPy} focuses on the latter while providing a reproducible framework for studying BMSSP's practical performance.
}

%\clearpage

\section{The BMSSPy Package}\label{sec:pack}

\texttt{BMSSPy}~\citep{makowski_bmsspy_2025} is {\color{blue}an open-source Python implementation} of the \emph{Bounded Multi-Source Shortest Path} (BMSSP) algorithm~\citep{duan_breaking_2025}. It supports directed graphs with nonnegative edge weights and returns both shortest-path \emph{distances} and \emph{predecessors} for path reconstruction. {\color{blue}Nodes are represented by integer indices \(0,\ldots,n-1\), and graphs are stored as adjacency lists. Installation and a complete usage example are provided in Appendix~\ref{app:install_example}.} Key features include:
\begin{itemize}
\item Support for either a single source or a set of source nodes.
\item Shortest-path distances and predecessor arrays for explicit path reconstruction.
\item Robust handling of unreachable nodes and variable node degrees.
\item Explicit implementation of theoretical components such as tight-edge forest construction, pivot selection, and the required data-structure operations.
\end{itemize}



{\color{blue}
\section{Functional Implementation of BMSSP}\label{sec:func}

This section describes the implementation of the Bounded Multi-Source Shortest Path (BMSSP) algorithm of \citet{duan_breaking_2025} in \texttt{BMSSPy}. At a high level, BMSSP avoids globally ordering all tentative distances, as Dijkstra does, by recursively processing nodes within bounded distance ranges. At each recursion level, Find Pivots (Algorithm~\ref{alg:bmssp_find_pivots}) reduces the active frontier, Recursive BMSSP (Algorithm~\ref{alg:bmssp_recursive_bmssp}) processes that frontier in bounded batches, and Base Case (Algorithm~\ref{alg:bmssp_base_case}) resolves sufficiently small subproblems using a bounded Dijkstra-style search. These three components compute shortest-path distances without globally sorting all nodes.

Our implementation retains these core components while making the additional steps needed for an executable end-to-end solver explicit. To distinguish these implementation additions and clarifications from the core BMSSP algorithm of \citet{duan_breaking_2025}, we use five color-coded categories throughout the pseudocode. The colors serve only as visual annotations of the type of implementation contribution; they do not represent algorithmic states or variables. Each colored modification is also identified by its corresponding acronym in the pseudocode, so the classification does not rely on color alone. Uncolored steps correspond to the underlying BMSSP procedure. The five categories are:

 }
\begin{enumerate}
\item \textcolor{Red}{\textbf{End-to-End Framework (E2E)}}:
We provide a complete workflow returning shortest path distances and predecessor arrays, supporting explicit path reconstruction for any given graph, origin set, and destination.

\item \textcolor{ForestGreen}{\textbf{Predecessor Tracking (PT)}}:
We explicitly show how to maintain a predecessor array, enabling the construction of full shortest paths beyond simple distance outputs.

\item \textcolor{Brown}{\textbf{Relaxation of Uniqueness Assumption (RUA)}}:
The original algorithm assumes that all path lengths are unique. The original paper offers a brief explanation that this {\color{blue}assumption does not lose generality after Lemma 2.1 in \cite{duan_breaking_2025}, and mentions one abstract way to implement this. We explicitly detail how to relax this uniqueness assumption to allow equal length paths} by using a distance value, a counter for nodes traversed, and a unique edge value system to provide lexicographic ordering.

\item \textcolor{Bittersweet}{\textbf{Explicit Forest and Pivot Construction (EFPC)}}:
We explicitly construct the tight-edge forest and derive the pivot set with rigorous criteria, clarifying previously implicit algorithmic steps.

\item \textcolor{Fuchsia}{\textbf{Fixes, Improvements, Notations (FIN)}}:
We implement several notation {\color{blue}and functional} clarifications that address theoretical and practical issues, ensuring reliable performance on arbitrary graphs. %TODO: should we mention our set work here?
\end{enumerate}

We first address the preliminary assumptions in Section~\ref{sec:handline_assumptions}. {\color{blue}Throughout the main algorithms in Section~\ref{subsec:alg_glance} and the end-to-end framework in Section~\ref{subsec:alg_framework}, each implementation addition is marked using the corresponding color and acronym defined above (\textcolor{Red}{E2E}, \textcolor{ForestGreen}{PT}, \textcolor{Brown}{RUA}, \textcolor{Bittersweet}{EFPC}, or \textcolor{Fuchsia}{FIN}). The underlying data structure is discussed in Appendix~\ref{app:data_structure}, and all notation follows Table~\ref{tab:notations}.}

 


\begin{table}[!htp]

\centering
\caption{Notation used throughout the paper.}
\label{tab:notations}

\begin{adjustbox}{
    max width=\textwidth,
    max totalheight=0.88\textheight,
    keepaspectratio,
    center
}
\begin{minipage}[t]{0.6\textwidth}
\scriptsize
\renewcommand{\arraystretch}{0.92}
\setlength{\tabcolsep}{3pt}

\begin{tabularx}{\linewidth}{@{}p{0.30\linewidth}X@{}}
\hline

\textbf{Notation} & \textbf{Description} \\
\hline

\multicolumn{2}{l}{\textbf{Graph Structure}} \\
\hline
$G=(V,E)$ & Directed graph with node set $V$ and edge set $E$. \\
$V$ & Set of nodes. \\
$E$ & Set of edges. \\
$n=|V|$ & Number of nodes. \\
$m=|E|$ & Number of edges. \\
$u,v\in V$ & Individual nodes. \\
$w$ & Nonnegative weight of edge $(u,v)$. \\
\hline

\multicolumn{2}{l}{\textbf{Distance, Paths, and Predecessors}} \\
\hline
$\text{dist}[v]$ & Tentative shortest-path distance to node $v$. \\
$\text{cdist}[v]$ & Counter-adjusted distance used for tie-breaking. \\
$\text{pred}[v]$ & Predecessor of $v$ in the shortest-path tree. \\
$\text{path}$ & Ordered list of nodes forming a shortest path. \\
$\text{length}$ & Total length of the returned path. \\
$\infty$ & Indicates an unreachable node. \\
$-1$ & Sentinel for no predecessor or an unreachable node. \\
\hline

\multicolumn{2}{l}{\textbf{Origins and Destinations}} \\
\hline
$O$ & Set of origin nodes. \\
$o$ & Origin node. \\
$\text{dest}$ & Destination node. \\
\hline

\multicolumn{2}{l}{\textbf{Distance Bounds}} \\
\hline
$B$ & Current global distance bound. \\
$B'$ or $B_{\text{new}}$ & Updated distance bound after refinement. \\
$B_{\text{curr}}$ & Active distance bound at the current recursion stage. \\
$B_{\text{comp}}$ & Completion bound for fully processed nodes. \\
\hline



\end{tabularx}
\end{minipage}
\hspace{0.015\textwidth}
\begin{minipage}[t]{0.6\textwidth}
\scriptsize
\renewcommand{\arraystretch}{0.92}
\setlength{\tabcolsep}{3pt}

\begin{tabularx}{\linewidth}{@{}p{0.30\linewidth}X@{}}
\hline
\textbf{Notation} & \textbf{Description} \\
\hline

\multicolumn{2}{l}{\textbf{Frontier and Recursion Sets}} \\
\hline
$S$ & Frontier set at the current recursion level. \\
$S'$ & New frontier set. \\
$S_{\mathrm{curr}}$ & Frontier subset currently being processed. \\
$U$ & Set of completed or processed nodes returned by a recursion. \\
$U_{\mathrm{curr}}$ & Nodes completed by the current recursive call. \\
$W$ & Expanded frontier produced during pivot finding. \\
$W_{\mathrm{prev}}$ & Frontier from the previous relaxation iteration. \\
$W_{\mathrm{curr}}$ & Frontier in the current relaxation iteration. \\
$P$ & Set of pivots selected for the next recursion level. \\
$F$ & Directed tight-edge forest over nodes in $W$. \\
$F^{ID}$ & Nodes with positive in-degree in $F$; forest roots are excluded. \\
\hline

\multicolumn{2}{l}{\textbf{Algorithm Parameters}} \\
\hline
$k=\lfloor(\log_2 n)^{1/3}\rfloor$
    & Frontier-growth, pivot, and base-case threshold. \\
$t=\lfloor(\log_2 n)^{2/3}\rfloor$
    & Controls work growth between recursion levels. \\
$l=\lceil\log_2(n)/t\rceil$
    & Recursion depth. \\
$M=2^{(l-1)t}$
    & Data-structure capacity at recursion level $l$. \\
$\text{precision}$
    & Number of decimal digits retained for edge weights. \\
\hline

\multicolumn{2}{l}{\textbf{Tie-Breaking and Structural Metadata}} \\
\hline
$\text{counter}$ & Tie-breaking increment used in effective distances. \\
$G_{\text{edge}}[u][v]$
    & Unique edge adjustment used for deterministic tie-breaking. \\
$\text{Heap}$ & Binary min-heap used for priority-queue operations. \\
\hline

\multicolumn{2}{l}{\textbf{Temporary and Data-Structure Variables}} \\
\hline
$D$ & Bounded-key data structure used by Recursive BMSSP. \\
$K$ & Temporary set of improved nodes passed to \textsc{BatchPrepend}. \\
$d'$ & Candidate effective distance produced by an edge relaxation. \\
$\mathit{edge\_adj\_value}$ & Scale used to assign unique edge adjustments. \\
$\mathit{edge\_id}$ & Unique adjustment associated with an edge. \\
\hline

\end{tabularx}
\end{minipage}
\end{adjustbox}

\end{table}

\subsection{Handling Assumptions}\label{sec:handline_assumptions}
The original paper assumes two properties of the input graph: (1) all path distances are unique, and (2) the graph has constant degree. The original paper also inherently assumes that common "set" operations happen in $O(1)$ time. In this section, we describe how we address these assumptions so that the algorithm can operate on graphs with nonnegative edge weights.

\subsubsection{Unique Distances}\label{sec:unique_distances}
The algorithm requires all path distances to be unique. In the original formulation, this uniqueness is ensured by augmenting each distance value with additional tie-breaking information, such as a counter or endpoint identifiers. Conceptually, one may represent each distance as a tuple
\[
\mathit{dist}[v] = (\mathit{length}, \mathit{path\_size}, \mathit{edge\_id}),
\]
where $\mathit{length}$ is the cumulative path weight, $\mathit{path\_size}$ records the number of edges used, and $\mathit{edge\_id}$ uniquely identifies the most recently added edge. During relaxation {\color{blue} of node u to node v}, this tuple updates via
\[
\mathit{dist}[v] \leftarrow \big(\mathit{dist}[u][0] + G[u][v], \mathit{dist}[u][1] + 1, G_{\mathit{edge}}[u][v]\big).
\]

This representation enables lexicographic comparison of distance values and guarantees uniqueness. The original paper proposes a slightly larger composite key:
\[
\mathit{dist}[v] = (\mathit{length},\; \mathit{path\_size},\; u,\; v),
\]
where the final two components serve as a canonical tie-breaking mechanism.

%While tuple-based lexicographic comparison is conceptually clean, it is operationally expensive: creating and accessing each tuple adds operational overhead, every comparison becomes a multi-field comparison, and auxiliary algorithms such as median-of-medians (used internally in \texttt{BMSSPy}) incur additional overhead and require algorithmic modifications. To avoid this, we encode the tuple as a \textit{single scalar value} while preserving the exact lexicographic order.

Tuple construction and direct comparisons would add substantial overhead to the repeated comparisons used throughout BMSSP. We therefore encode the same lexicographic ordering in a single scalar value by placing the path-size and edge-id terms below the user-specified decimal precision.

We first round $\mathit{length}$ to a fixed number of digits provided as $\mathit{precision}$. Then we {\color{blue}scale the order of magnitude of} the decimal representation of $\mathit{path\_size}$ and $\mathit{edge\_id}$ sufficiently so that the three components never overlap. Let
\begin{equation*}
\begin{split}
\mathit{counter} &= 10^{-(\mathit{precision} + \ceil{\log_{10}(2n + 1)})} \\
\mathit{edge\_adj\_value} &= 10^{-(\mathit{precision} + \ceil{\log_{10}(2n + 1)} + \ceil{\log_{10}(m + 1)})}
\end{split}
\end{equation*}
We use $2n+1$ to allow simple rounding of the distance to the precision level. If instead we chose to truncate the distance to that precision level, we could use $n+1$. Continuing, the scalar distance is then
\begin{equation*}
\begin{split}
\mathit{dist}[v] =\;& \operatorname{round}(\mathit{length}, \mathit{precision}) + \mathit{path\_size} \cdot \mathit{counter} + \mathit{edge\_id} \cdot \mathit{edge\_adj\_value}.
\end{split}
\end{equation*}

Because each edge has a fixed unique identifier, $G_{\mathit{edge}}[\cdot][\cdot]$ is preprocessed so that it already stores the values $\{\mathit{edge\_adj\_value}, 2\,\mathit{edge\_adj\_value}, \dots, m\,\mathit{edge\_adj\_value}\}$. During relaxation, we update
\begin{equation*}
\begin{split}
\mathit{cdist}[v] &\leftarrow \mathit{cdist}[u] + G[u][v] + \mathit{counter} \\
\mathit{dist}[v] &\leftarrow \mathit{cdist}[u] + G[u][v] + \mathit{counter} + G_{\mathit{edge}}[u][v],
\end{split}
\end{equation*}
where $\mathit{cdist}$ omits the edge id portion, allowing us to skip rounding operations at each relaxation.

Because this approach relies on large decimal shifts, it is susceptible to floating-point error. Therefore, the implementation uses Python's \texttt{Decimal} type to preserve the small deterministic tie-breaking offsets at the specified precision.

\subsubsection{Constant-Degree Transformation}
The stated BMSSP complexity bound assumes a constant-degree representation. For graphs that do not satisfy this condition, \texttt{BMSSPy} can first transform the graph into the constant-degree representation described in Algorithm~\ref{alg:utils_convert_to_constant_degree} in Appendix~\ref{app:to_out_degree}. Auxiliary zero-weight edges preserve path lengths, and Algorithm~\ref{alg:utils_convert_from_constant_degree} in Appendix~\ref{app:from_out_degree} maps the resulting distances and predecessors back to the original graph. We refer to experiments using this transformation as BMSSP-CD; BMSSP without the transformation is reported separately.

}

\subsubsection{Set and Dictionary Operations}\label{subsec:set_functions}

{\color{blue}
BMSSP repeatedly requires set membership, insertion, and key lookup in deterministic \(O(1)\) time. Standard Python sets and dictionaries are hash based and therefore provide expected, rather than deterministic worst-case, constant-time access. Because nodes in \texttt{BMSSPy} are represented by integer indices \(0,\ldots,n-1\), we instead use an index based membership array together with a list of included set values and a generation counter to validate current state. This provides deterministic \(O(1)\) membership and insertion while retaining \(O(|S|)\) iteration and set extension. The same principle is extended with direct-indexed value arrays to provide deterministic dictionary-style lookup, assignment, and invalidation. 

Independent state needs to be maintained across recursion levels because each level of \textsc{RecursiveBMSSP} has its own sets and data-structure state. Since the recursion depth is \(l=\Theta((\log n)^{1/3})\), initializing these sets one time requires \(O(n(\log n)^{1/3})\) time and storage. The implementation mechanics and their integration with the BMSSP data structure are described in Appendix~\ref{app:set_dictionary_details} and Appendix~\ref{app:data_structure}.

}

% \clearpage
\subsection{Main Algorithms}\label{subsec:alg_glance}

{\color{blue}
BMSSP is organized around three recursive components, corresponding to Algorithms~1--3 of \citet{duan_breaking_2025}. Their roles can be understood at a high level as follows. Find Pivots (Algorithm~\ref{alg:bmssp_find_pivots}) reduces a potentially large frontier to a smaller set of representative nodes. Recursive BMSSP (Algorithm~\ref{alg:bmssp_recursive_bmssp}) repeatedly processes these representatives within bounded distance windows and passes smaller subproblems down the recursion. When the recursion reaches \(l=0\), Base Case (Algorithm~\ref{alg:bmssp_base_case}) completes a bounded Dijkstra-style search. The combination replaces Dijkstra's single global priority ordering with a hierarchy of locally bounded searches.
}


\paragraph{Algorithm~\ref{alg:bmssp_find_pivots}: Find Pivots.}
This subroutine expands the current frontier $S$ under a distance bound $B$ and selects a reduced subset of \emph{pivots} $P \subseteq S$ for the next recursion level. It performs $k$ rounds of bounded relaxations (Bellman--Ford style) to build a temporary frontier $W$ of nodes with distance $< B$ from $S$. If $W$ grows too quickly ($|W| > k|S|$), the routine short-circuits and sets $P = S$, capping the work passed downward. Otherwise, it constructs a \emph{tight-edge forest} on $W$ (via the predecessor relation) and chooses as pivots those $u \in S$ that are forest roots with subtree size at least $k$. These pivots summarize the influence of $S$ in the current window, reducing redundant relaxations and avoiding global ordering while preserving the theoretical recursion and batch bounds.


Our Find Pivots implementation in \texttt{BMSSPy} refines several steps of the original theoretical description to ensure correctness and clarity on arbitrary weighted graphs:
\begin{itemize} 
    \item \textit{Use of evolving frontier variables} (\text{Lines~\ref{alg:bmssp_find_pivots_start}--\ref{alg:bmssp_find_pivots_prev_update}}): We replace indexed frontier arrays ($W_i$) with $W_{\mathrm{curr}}$ and $W_{\mathrm{prev}}$ for clarity and memory efficiency. 
    \item \textit{Counter-based unique distances} (\text{Lines~\ref{alg:bmssp_find_pivots_counter_dist_update} and ~\ref{alg:bmssp_find_pivots_counter_dist_update2}}): We apply our lightweight tie-breaking mechanism described in Section \ref{sec:unique_distances} using $counter$ and $G_{\text{edge}}$ adjustments so that all effective distances remain unique. The relevant parameters are initialized in Algorithm~\ref{alg:bmssp_wrapper_init}. 
    \item \textit{Predecessor tracking} (\text{Line~\ref{alg:bmssp_find_pivots_pred}}): updates a global \texttt{pred[]} array defined in Algorithm~\ref{alg:bmssp_solve} for path reconstruction and for later pivot determination in the tight-edge forest. 
    \item \textit{Explicit forest construction} (\text{Lines~\ref{alg:bmssp_find_pivots_forest_init}--\ref{alg:bmssp_find_pivots_forest_loop_end}}): builds the adjacency-list representation of the tight-edge forest $F$ directly from recorded predecessors rather than distance equalities. 
    \item \textit{Pivot root selection} (\text{Lines~\ref{alg:bmssp_find_pivots_roots_start}--\ref{alg:bmssp_find_pivots_roots_end}}): explicitly checks in-degree $u \notin F^{ID}$ to identify forest roots and calls \textsc{IsPivot} ({\color{blue}Appendix~\ref{app:is_pivot} - Algorithm~\ref{alg:utils_is_pivot}}) to determine whether their subtrees exceed the size threshold $k$. 
\end{itemize}

\input{formulas/bmssp/find_pivots}

\paragraph{Algorithm~\ref{alg:bmssp_base_case}: Base Case.}
At recursion depth $l=0$, BMSSP reduces to a bounded single-source search from the unique {\color{blue}node} $x\in S$. We run a Dijkstra-style exploration under the current bound $B$, discovering nodes with distance $<B$ whose shortest paths visit $x$. The search stops as soon as the $(k+1)$st node would be explored; its weight becomes the refined boundary $B'$, and the already-settled $k$ nodes form the new frontier below $B'$. If fewer than $k+1$ nodes are found, we return $(B',S')=(B,S')$. Our Base Case implementation in \texttt{BMSSPy} makes the following clarifications to ensure correctness:
\begin{itemize}
  \item \textit{Early termination at $k{+}1$} (Line~\ref{alg:bmssp_base_case_early_break}): set $B'$ to the weight of the would-be $(k{+}1)$st node and stop, yielding exactly $k$ completed nodes below $B'$. This is done slightly differently than the original \citet{duan_breaking_2025} Algorithm 2 in that we avoid the extra operational overhead of filtering for only values lower than $B'$ after the main while loop ends.
  \item \textit{Counter-based unique distances} (Lines~\ref{alg:bmssp_base_case_modified_distance} and \ref{alg:bmssp_base_case_counter_dist_update}): similar to Algorithm \ref{alg:bmssp_find_pivots}.
  \item \textit{Predecessor tracking} (Line~\ref{alg:bmssp_base_case_pred}): update \texttt{pred[]} on improvement, as in Algorithm \ref{alg:bmssp_find_pivots}.
\end{itemize}

\input{formulas/bmssp/base_case}

\paragraph{Algorithm~\ref{alg:bmssp_recursive_bmssp}: Recursive BMSSP.}
For $l>0$, this routine coordinates the BMSSP recursion. The algorithm first compresses the frontier using \textsc{FindPivots} to obtain the pivot set $P$ and temporary frontier $W$. It then initializes the bounded data structure $D$ with capacity $M=2^{(l-1)t}$ and inserts the pivots into it. While the work budget is not exceeded and the frontier remains nonempty, it:
(i) pulls the next bounded batch $(B_{\mathrm{curr}},S_{\mathrm{curr}})$,
(ii) recursively processes this batch at level $l-1$,
(iii) relaxes outgoing edges from the newly completed set $U_{\mathrm{curr}}$, and
(iv) returns improved nodes to the appropriate distance window, either $[B_{\mathrm{curr}},B)$ through \textsc{Insert} or $[B_{\mathrm{comp}},B_{\mathrm{curr}})$ through \textsc{BatchPrepend}.
It returns the completion bound $B_{\mathrm{comp}}$ together with the nodes completed below that bound.

Because the deterministic set representation described in Section~\ref{subsec:set_functions} stores node identifiers rather than key--value tuples, the temporary set $K$ stores only node identifiers (Line~\ref{alg:bmssp_recursive_collectK}). Their current distances are attached when \textsc{BatchPrepend} is called (Line~\ref{alg:bmssp_recursive_batchprepend}). Recursive BMSSP also applies the counter-based effective-distance construction of Section~\ref{sec:unique_distances} when computing candidate keys and updating counter distances (Lines~\ref{alg:bmssp_recursive_key} and~\ref{alg:bmssp_recursive_key2}), and updates the predecessor array whenever a distance improves (Line~\ref{alg:bmssp_recursive_bmssp_pred}), consistent with Algorithms~\ref{alg:bmssp_find_pivots} and~\ref{alg:bmssp_base_case}.

\input{formulas/bmssp/recursive_bmssp}


\subsection{End-to-End Framework}\label{subsec:alg_framework}

In the previous section, the BMSSP framework comprises three main algorithms: Find Pivots (Algorithm~\ref{alg:bmssp_find_pivots}), Base Case (Algorithm~\ref{alg:bmssp_base_case}), and Recursive BMSSP (Algorithm~\ref{alg:bmssp_recursive_bmssp}), where the latter coordinates the first two.

Our main contribution is an end-to-end implementation that utilizes BMSSP to compute shortest-path distances and reconstruct predecessor paths for a given graph~$G$. As illustrated in Figure~\ref{fig:framework}, the workflow proceeds as follows:

\begin{enumerate}
  \item \textbf{Initialization} - Algorithm~\ref{alg:bmssp_wrapper_init} preprocesses the input graph $G$ (in adjacency-list form) and sets up the precision parameter~$precision$, along with adjustment values (\textit{counter}, $G_{\text{edge}}$) to ensure unique path distances.
  \item \textbf{Solve phase} - Algorithm~\ref{alg:bmssp_solve} initializes distance and predecessor arrays, sets parameters $k$, $t$, and $l$, and invokes the Recursive BMSSP (Algorithm~\ref{alg:bmssp_recursive_bmssp}).
  \item \textbf{Post-processing} - Distances are rounded according to the specified precision level, and final paths are reconstructed using Algorithm~\ref{alg:utils_reconstruct_path} in Appendix \ref{app:path} (Algorithm \ref{alg:utils_convert_from_constant_degree} in Appendix \ref{app:from_out_degree} is used prior to Algorithm~\ref{alg:utils_reconstruct_path} if the graph is converted to constant degree).
\end{enumerate}

This end-to-end process returns both shortest-path distances and predecessor paths for the given origin set~$O$ and, optionally, a destination node~$dest$. The following sections describe the remaining algorithms that complete the framework.

\begin{figure}[!ht]
\begin{center}
\caption{End-to-End BMSSPy Framework}\label{fig:framework}
\includegraphics[width=\linewidth]{figures/framework.pdf}
\end{center}
\end{figure}


\paragraph{Algorithm~\ref{alg:bmssp_solve}: BMSSP Solve.}
This is the main entry point of the framework. It initializes the distance, counter-distance, and predecessor arrays, sets algorithm parameters 
($k$, $t$, and $l$ using $\log_2$), and initializes structures for sets and dictionaries here (see Section \ref{subsec:set_functions} and Appendix~\ref{app:data_structure}). We then invoke the Recursive BMSSP (Algorithm~\ref{alg:bmssp_recursive_bmssp}). After computing distances, 
it rounds them back to the specified precision and optionally reconstructs the shortest path to a target node using Algorithm~\ref{alg:utils_reconstruct_path} in {\color{blue}Appendix~\ref{app:path}}. 
Our implementation also clarifies the $\log_2$ base used in parameter definitions and integrates explicit path reconstruction, returning both distance and path results.
\input{formulas/bmssp/solve}

\paragraph{Algorithm~\ref{alg:bmssp_wrapper_init}: BMSSP Initialization.}
This setup routine preprocesses the input graph and prepares adjustment constants needed to meet the unique path length requirements for the algorithm. 
It rounds edge weights to the desired precision~$p$, computes unique offset parameters (\texttt{counter} and $G_{\text{edge}}$) to ensure unique distances for paths, 
and stores these within the BMSSP object. The resulting preprocessed graph supports deterministic, precision-consistent shortest-path calculations.

\input{formulas/bmssp/initialize}




%\clearpage
{\color{blue}
\section{Numerical Experiments}\label{sec:num}
We evaluate the practical performance of BMSSP using both Python and C++ implementations. The experiments address three questions: (i) how does BMSSP compare with Dijkstra implementations on synthetic gridgraphs and real geospatial networks (Section~\ref{subsec:compar}); (ii) under extrapolation, at what graph size, if any, might BMSSP overtake \texttt{SCGraph} Dijkstra (Section~\ref{subsec:break}); and (iii) how sensitive are BMSSPy's accuracy and runtime to its implementation-specific precision parameter? (Section~\ref{subsec:precision_sensitivity}).



\subsection{Comparison of Run Times on Grid Graphs and Geospatial Graphs}\label{subsec:compar}
We compare BMSSP in Python and C++, both with and without the constant-degree transformation (BMSSP-CD), against five Dijkstra implementations from four baseline families: \emph{Vanilla Dijkstra}, a simple pure-Python implementation without a heap; \texttt{SCGraph} Dijkstra, evaluated in both Python and C++; \texttt{NetworkX} Dijkstra; and \texttt{igraph} Dijkstra. For each package baseline, we use its standard Dijkstra implementation. We exclude $A^*$ ~\citep{hart1968formal} because \texttt{SCGraph} Dijkstra outperforms \texttt{SCGraph} $A^*$ when shortest paths deviate substantially from straight-line geometry \citep{makowski2025scgraph}.

\paragraph{Setup:}
We evaluate two families of synthetic gridgraphs: square and rectangular. Each contains a vertical obstacle extending from the top center to five grid units above the bottom boundary. Square instances range from \(15\times15\) to \(1000\times1000\), and rectangular instances range from \(15\times50\) to \(100\times2000\). Each instance is evaluated from three origins: Bottom Left (BL), Top Right (TR), and Center (C). Figure~\ref{fig:shortest_path_trees} illustrates the shortest-path trees of BMSSP-CD for the three origins on a \(50\times50\) square gridgraph \citep{makowski2025scgraph}.
We also evaluate five real geospatial networks from \texttt{SCGraph} and \texttt{SCGraph\_Data}~\citep{makowski2025scgraph}: Marnet, North America Rail, Oak Ridge Maritime, US Freeway, and the combined World Highways and Marnet network. Each is evaluated from Los Angeles (LA), New York (NY), and Seattle (SEA). Figure~\ref{fig:na_railway} shows an example shortest-path tree from Los Angeles on the North America Rail network.


All experiments were executed on Ubuntu~24.04.3, running single-threaded on an AMD Ryzen™ Threadripper™ 3970X 32-core 4.5 GHz machine with 128 GiB RAM. The Python implementation uses Python 3.12.3. The C++ implementation was compiled against the C++20 standard using gcc 13.3.0 with the flags: ``-O3 -DNDEBUG". Both implementations follow the pseudocode presented in this paper.

\begin{figure}[!ht]
\begin{center}
\caption{Example Construction of Shortest-path Tree for 50×50 Square Gridgraph for Bottom Left (BL), Top Right (TR), and Center (C) Cases}\label{fig:shortest_path_trees}
\includegraphics[width=\linewidth]{figures/shortest_path_trees1.pdf}
\end{center}
\end{figure}

\begin{figure}[!ht]
\begin{center}
\caption{Example Construction of Shortest-path Tree from Los Angeles for North American Rail}\label{fig:na_railway}
\includegraphics[width=0.8\linewidth]{figures/north_american.png}
\end{center}
\end{figure}


\paragraph{Results.}
Figure~\ref{fig:algorithm_runtimes} summarizes the computational results in
absolute runtime and relative to both the Python and C++ implementations of
BMSSP. Tables~\ref{tab:square_gridgraphs}, \ref{tab:rect_gridgraphs}, and
\ref{tab:geospatial} report the corresponding mean runtimes and standard
deviations for each graph combination. 
\begin{figure}[!ht]
\begin{center}
\caption{
Runtime comparison across square gridgraphs, rectangular gridgraphs, and
geospatial graphs.
Panels (a)--(c) report average absolute runtimes (original scale and bands in Figure \ref{fig:algorithm_runtimes_app} in Appendix \ref{app:runtime_discussion});
Panels (d)--(f) report runtimes relative to Python BMSSP; and
Panels (g)--(i) report runtimes relative to BMSSP C++.
A relative multiplier greater than one indicates that an implementation is
slower than the corresponding BMSSP reference, whereas a value below one
indicates that it is faster.
}
\label{fig:algorithm_runtimes}
\includegraphics[width=\linewidth]{figures/algorithm_runtimes3x3.pdf}
\end{center}
\end{figure}


\begin{table}[!htp]
\caption{\textbf{Runtime Comparison on Square Gridgraphs
(Mean $\pm$ SD over 10 runs; ms)}}
\label{tab:square_gridgraphs}
\centering
\resizebox{\linewidth}{!}{%
\scriptsize\color{blue}
\begin{tabular}{
    l l r r
    r@{$\,\pm\,$}l
    r@{$\,\pm\,$}l
    r@{$\,\pm\,$}l
    r@{$\,\pm\,$}l
    r@{$\,\pm\,$}l
    r@{$\,\pm\,$}l
    r@{$\,\pm\,$}l
    r@{$\,\pm\,$}l
    r@{$\,\pm\,$}l
}
\toprule
\textbf{\shortstack[c]{Grid Graph\\Square}} &
\raisebox{0.5\normalbaselineskip}{\textbf{Case}} &
\raisebox{0.5\normalbaselineskip}{\textbf{Nodes}} &
\raisebox{0.5\normalbaselineskip}{\textbf{Edges}} &
\multicolumn{2}{c}{\raisebox{0.5\normalbaselineskip}{\textbf{BMSSP-CD}}} &
\multicolumn{2}{c}{\textbf{\shortstack[c]{BMSSP-CD\\(C++)}}} &
\multicolumn{2}{c}{\raisebox{0.5\normalbaselineskip}{\textbf{BMSSP}}} &
\multicolumn{2}{c}{\textbf{\shortstack[c]{BMSSP\\(C++)}}} &
\multicolumn{2}{c}{\textbf{\shortstack[c]{Vanilla\\Dijkstra}}} &
\multicolumn{2}{c}{\textbf{\shortstack[c]{SCGraph\\Dijkstra}}} &
\multicolumn{2}{c}{\textbf{\shortstack[c]{SCGraph Dijkstra\\(C++)}}} &
\multicolumn{2}{c}{\textbf{\shortstack[c]{NetworkX\\Dijkstra}}} &
\multicolumn{2}{c}{\textbf{\shortstack[c]{igraph\\Dijkstra}}} \\
\midrule

15$\times$15 & BL & 225 & 1,484 &
37.61 & 0.23 & 2.82 & 0.10 & 7.92 & 0.05 & 1.06 & 0.02 &
0.83 & 0.01 & 0.17 & 0.01 & 0.09 & 0.22 & 0.42 & 0.02 & 0.11 & 0.03 \\
15$\times$15 & TR & 225 & 1,484 &
36.73 & 0.21 & 2.73 & 0.11 & 7.71 & 0.06 & 1.01 & 0.01 &
0.82 & 0.01 & 0.16 & 0.01 & 0.02 & 0.02 & 0.42 & 0.02 & 0.12 & 0.02 \\
15$\times$15 & C & 225 & 1,484 &
38.18 & 0.21 & 2.84 & 0.10 & 7.84 & 0.12 & 1.03 & 0.02 &
0.83 & 0.01 & 0.16 & 0.01 & 0.02 & 0.02 & 0.42 & 0.02 & 0.11 & 0.02 \\
\midrule

25$\times$25 & BL & 625 & 4,424 &
113.99 & 0.45 & 8.39 & 0.26 & 11.39 & 0.14 & 1.58 & 0.03 &
5.32 & 0.01 & 0.52 & 0.01 & 0.08 & 0.06 & 1.41 & 0.10 & 0.41 & 0.05 \\
25$\times$25 & TR & 625 & 4,424 &
89.78 & 0.51 & 6.56 & 0.32 & 13.85 & 0.20 & 1.97 & 0.03 &
5.37 & 0.06 & 0.47 & 0.01 & 0.07 & 0.07 & 1.39 & 0.10 & 0.49 & 0.07 \\
25$\times$25 & C & 625 & 4,424 &
107.21 & 0.44 & 7.86 & 0.24 & 11.59 & 0.16 & 1.64 & 0.03 &
5.38 & 0.02 & 0.49 & 0.01 & 0.07 & 0.06 & 1.42 & 0.09 & 0.42 & 0.05 \\
\midrule

50$\times$50 & BL & 2,500 & 18,774 &
540.34 & 1.67 & 39.06 & 0.33 & 38.23 & 0.48 & 5.72 & 0.09 &
78.08 & 0.08 & 1.97 & 0.01 & 0.28 & 0.13 & 6.04 & 0.36 & 3.07 & 0.21 \\
50$\times$50 & TR & 2,500 & 18,774 &
446.44 & 1.33 & 31.98 & 0.46 & 56.81 & 0.50 & 8.08 & 0.09 &
78.49 & 0.12 & 1.92 & 0.01 & 0.28 & 0.13 & 6.12 & 0.40 & 3.80 & 0.27 \\
50$\times$50 & C & 2,500 & 18,774 &
518.83 & 1.35 & 37.50 & 0.39 & 41.42 & 0.63 & 6.06 & 0.09 &
78.63 & 0.12 & 1.97 & 0.01 & 0.30 & 0.13 & 6.11 & 0.34 & 2.97 & 0.22 \\
\midrule

75$\times$75 & BL & 5,625 & 43,124 &
1,300.02 & 250.47 & 83.55 & 0.17 & 82.26 & 0.34 & 12.75 & 0.12 &
403.02 & 0.62 & 4.36 & 0.02 & 0.88 & 1.00 & 15.22 & 0.61 & 10.99 & 0.44 \\
75$\times$75 & TR & 5,625 & 43,124 &
1,038.59 & 247.56 & 74.32 & 0.30 & 98.64 & 0.45 & 13.98 & 0.06 &
402.88 & 0.62 & 4.36 & 0.03 & 0.67 & 0.48 & 16.11 & 0.78 & 19.53 & 0.47 \\
75$\times$75 & C & 5,625 & 43,124 &
1,185.09 & 257.44 & 76.69 & 0.23 & 88.26 & 0.53 & 13.03 & 0.10 &
403.09 & 0.85 & 4.49 & 0.02 & 0.70 & 0.41 & 15.38 & 0.65 & 11.45 & 0.64 \\
\midrule

100$\times$100 & BL & 10,000 & 77,474 &
2,454.75 & 273.58 & 161.25 & 0.26 & 152.90 & 0.25 & 23.55 & 0.07 &
1,353.53 & 7.55 & 7.65 & 0.06 & 1.60 & 1.80 & 48.08 & 46.37 & 35.58 & 0.36 \\
100$\times$100 & TR & 10,000 & 77,474 &
2,028.55 & 358.08 & 167.79 & 0.34 & 171.70 & 0.43 & 25.80 & 0.22 &
1,336.40 & 20.16 & 7.62 & 0.25 & 1.30 & 1.22 & 50.86 & 46.41 & 45.35 & 0.19 \\
100$\times$100 & C & 10,000 & 77,474 &
2,365.24 & 347.05 & 149.95 & 0.26 & 158.05 & 0.78 & 24.74 & 0.53 &
1,323.14 & 1.56 & 8.01 & 0.21 & 1.37 & 1.11 & 48.74 & 47.48 & 33.16 & 0.36 \\
\midrule

250$\times$250 & BL & 62,500 & 493,574 &
17,832.88 & 607.81 & 1,196.58 & 0.71 & 899.55 & 2.38 & 147.61 & 0.22 &
57,950.78 & 848.54 & 53.23 & 0.39 & 10.18 & 3.05 & 469.84 & 114.43 & 599.20 & 5.76 \\
250$\times$250 & TR & 62,500 & 493,574 &
13,736.76 & 643.98 & 871.22 & 0.94 & 1,072.57 & 72.21 & 165.27 & 0.14 &
56,140.88 & 578.25 & 64.88 & 0.66 & 15.16 & 2.81 & 488.90 & 118.37 & 757.55 & 5.90 \\
250$\times$250 & C & 62,500 & 493,574 &
16,732.32 & 642.33 & 1,127.00 & 1.11 & 1,007.96 & 69.00 & 157.61 & 0.41 &
56,714.30 & 575.72 & 58.88 & 0.55 & 11.41 & 3.09 & 426.37 & 106.84 & 553.02 & 57.02 \\
\midrule

500$\times$500 & BL & 250,000 & 1,987,074 &
73,607.67 & 1,756.16 & 4,577.18 & 52.87 & 3,746.24 & 133.60 & 606.04 & 3.12 &
\multicolumn{2}{c}{N/A} & 248.66 & 0.55 & 39.68 & 2.88 & 3,021.86 & 269.52 & 4,951.56 & 86.97 \\
500$\times$500 & TR & 250,000 & 1,987,074 &
57,275.22 & 1,592.46 & 3,571.67 & 15.62 & 4,393.82 & 125.46 & 687.66 & 14.80 &
\multicolumn{2}{c}{N/A} & 267.69 & 0.36 & 59.15 & 2.78 & 3,124.22 & 273.12 & 6,972.39 & 31.20 \\
500$\times$500 & C & 250,000 & 1,987,074 &
69,502.48 & 1,581.21 & 4,577.42 & 54.74 & 4,100.77 & 170.12 & 649.33 & 2.97 &
\multicolumn{2}{c}{N/A} & 275.67 & 1.17 & 51.10 & 2.66 & 2,720.93 & 252.34 & 4,613.97 & 100.16 \\
\midrule

750$\times$750 & BL & 562,500 & 4,480,574 &
170,943.34 & 719.78 & 10,145.74 & 121.32 & 8,147.52 & 237.54 & 1,387.79 & 20.41 &
\multicolumn{2}{c}{N/A} & 596.75 & 1.79 & 104.87 & 4.48 & 11,495.89 & 853.16 & 16,591.28 & 181.74 \\
750$\times$750 & TR & 562,500 & 4,480,574 &
137,793.11 & 2,655.11 & 8,406.71 & 92.83 & 9,648.46 & 210.56 & 1,586.22 & 7.66 &
\multicolumn{2}{c}{N/A} & 661.56 & 0.85 & 135.87 & 2.88 & 11,232.98 & 780.63 & 26,196.78 & 130.70 \\
750$\times$750 & C & 562,500 & 4,480,574 &
165,398.71 & 3,545.50 & 9,977.21 & 9.34 & 9,000.94 & 309.02 & 1,498.70 & 3.66 &
\multicolumn{2}{c}{N/A} & 623.15 & 2.27 & 125.10 & 3.26 & 9,689.92 & 693.53 & 15,497.07 & 164.85 \\
\midrule

1000$\times$1000 & BL & 1,000,000 & 7,974,074 &
323,176.31 & 4,989.07 & 19,291.17 & 147.61 & 14,925.54 & 470.21 & 2,522.75 & 22.31 &
\multicolumn{2}{c}{N/A} & 1,041.37 & 3.10 & 159.90 & 9.04 & 31,843.25 & 3,360.87 & 42,081.52 & 803.65 \\
1000$\times$1000 & TR & 1,000,000 & 7,974,074 &
242,573.89 & 4,039.83 & 15,135.21 & 171.49 & 17,447.94 & 385.70 & 3,053.33 & 13.88 &
\multicolumn{2}{c}{N/A} & 1,278.06 & 3.74 & 262.92 & 4.91 & 30,890.19 & 2,440.42 & 74,034.14 & 240.44 \\
1000$\times$1000 & C & 1,000,000 & 7,974,074 &
299,463.90 & 5,300.32 & 18,407.42 & 20.59 & 16,315.89 & 417.96 & 2,899.94 & 17.28 &
\multicolumn{2}{c}{N/A} & 1,128.19 & 8.51 & 218.60 & 5.66 & 25,410.46 & 2,179.69 & 38,717.81 & 253.56 \\

\bottomrule
\end{tabular}
}
%\vspace{2pt}
\parbox{\linewidth}{\tiny
\color{blue}
\textit{Notes:} BL = bottom-left origin; TR = top-right origin;
C = center origin. N/A indicates that Vanilla Dijkstra was not
evaluated for that instance due to extensive runtime.
}
\end{table}

\begin{table}[!htp]
\caption{\textbf{Runtime Comparison on Rectangular Gridgraphs
(Mean $\pm$ SD over 10 runs; ms)}}
\label{tab:rect_gridgraphs}
\centering
\resizebox{\linewidth}{!}{%
\scriptsize\color{blue}
\begin{tabular}{
    l l r r
    r@{$\,\pm\,$}l
    r@{$\,\pm\,$}l
    r@{$\,\pm\,$}l
    r@{$\,\pm\,$}l
    r@{$\,\pm\,$}l
    r@{$\,\pm\,$}l
    r@{$\,\pm\,$}l
    r@{$\,\pm\,$}l
    r@{$\,\pm\,$}l
}
\toprule
\textbf{\shortstack[c]{Grid Graph\\Rectangular}} &
\raisebox{0.5\normalbaselineskip}{\textbf{Case}} &
\raisebox{0.5\normalbaselineskip}{\textbf{Nodes}} &
\raisebox{0.5\normalbaselineskip}{\textbf{Edges}} &
\multicolumn{2}{c}{\raisebox{0.5\normalbaselineskip}{\textbf{BMSSP-CD}}} &
\multicolumn{2}{c}{\textbf{\shortstack[c]{BMSSP-CD\\(C++)}}} &
\multicolumn{2}{c}{\raisebox{0.5\normalbaselineskip}{\textbf{BMSSP}}} &
\multicolumn{2}{c}{\textbf{\shortstack[c]{BMSSP\\(C++)}}} &
\multicolumn{2}{c}{\textbf{\shortstack[c]{Vanilla\\Dijkstra}}} &
\multicolumn{2}{c}{\textbf{\shortstack[c]{SCGraph\\Dijkstra}}} &
\multicolumn{2}{c}{\textbf{\shortstack[c]{SCGraph Dijkstra\\(C++)}}} &
\multicolumn{2}{c}{\textbf{\shortstack[c]{NetworkX\\Dijkstra}}} &
\multicolumn{2}{c}{\textbf{\shortstack[c]{igraph\\Dijkstra}}} \\
\midrule

15$\times$50 & BL & 750 & 4,984 &
99.46 & 0.40 & 7.07 & 0.25 & 12.67 & 0.13 & 1.96 & 0.03 &
7.30 & 0.04 & 0.62 & 0.07 & 0.17 & 0.32 & 1.61 & 0.16 & 0.58 & 0.07 \\
15$\times$50 & TR & 750 & 4,984 &
96.21 & 0.40 & 6.49 & 0.25 & 22.43 & 0.14 & 2.81 & 0.04 &
7.38 & 0.12 & 0.53 & 0.01 & 0.07 & 0.07 & 1.61 & 0.12 & 0.89 & 0.11 \\
15$\times$50 & C & 750 & 4,984 &
92.28 & 0.40 & 6.47 & 0.29 & 16.77 & 0.21 & 2.45 & 0.14 &
7.41 & 0.12 & 0.53 & 0.01 & 0.08 & 0.10 & 1.61 & 0.11 & 0.67 & 0.08 \\
\midrule

25$\times$50 & BL & 1,250 & 8,924 &
234.80 & 0.58 & 16.95 & 0.35 & 19.63 & 0.22 & 3.09 & 0.04 &
19.88 & 0.03 & 1.06 & 0.01 & 0.14 & 0.06 & 2.91 & 0.15 & 1.23 & 0.11 \\
25$\times$50 & TR & 1,250 & 8,924 &
164.42 & 0.51 & 11.60 & 0.43 & 30.51 & 0.25 & 4.78 & 0.46 &
20.01 & 0.05 & 0.94 & 0.01 & 0.12 & 0.06 & 2.92 & 0.18 & 1.78 & 0.19 \\
25$\times$50 & C & 1,250 & 8,924 &
200.93 & 0.60 & 14.25 & 0.37 & 25.18 & 0.43 & 3.71 & 0.12 &
19.97 & 0.03 & 0.95 & 0.01 & 0.13 & 0.07 & 2.88 & 0.17 & 1.35 & 0.14 \\
\midrule

25$\times$100 & BL & 2,500 & 17,924 &
470.47 & 1.39 & 33.83 & 0.35 & 42.76 & 0.40 & 6.81 & 0.15 &
76.76 & 0.08 & 1.98 & 0.01 & 0.29 & 0.15 & 5.97 & 0.32 & 3.91 & 0.37 \\
25$\times$100 & TR & 2,500 & 17,924 &
308.16 & 1.33 & 25.73 & 0.49 & 83.54 & 0.62 & 10.81 & 0.67 &
76.78 & 0.20 & 1.82 & 0.01 & 0.25 & 0.13 & 6.26 & 0.41 & 7.05 & 0.73 \\
25$\times$100 & C & 2,500 & 17,924 &
387.55 & 1.01 & 27.20 & 0.36 & 71.31 & 0.75 & 9.27 & 0.21 &
77.38 & 0.11 & 1.91 & 0.01 & 0.27 & 0.13 & 6.12 & 0.38 & 4.67 & 0.46 \\
\midrule

50$\times$100 & BL & 5,000 & 37,774 &
985.29 & 3.45 & 67.48 & 0.20 & 73.07 & 0.49 & 12.27 & 0.52 &
315.63 & 0.85 & 3.84 & 0.03 & 0.55 & 0.32 & 13.29 & 0.63 & 11.16 & 0.74 \\
50$\times$100 & TR & 5,000 & 37,774 &
689.83 & 2.81 & 59.50 & 0.29 & 119.27 & 0.55 & 18.62 & 0.98 &
311.49 & 0.74 & 3.68 & 0.03 & 0.73 & 0.94 & 14.09 & 0.82 & 21.61 & 0.49 \\
50$\times$100 & C & 5,000 & 37,774 &
825.28 & 1.79 & 55.85 & 0.28 & 79.51 & 0.67 & 12.05 & 0.28 &
316.54 & 6.14 & 3.88 & 0.02 & 0.60 & 0.39 & 115.78 & 306.40 & 12.85 & 0.79 \\
\midrule

50$\times$200 & BL & 10,000 & 75,774 &
2,074.26 & 6.07 & 140.02 & 0.40 & 141.58 & 0.57 & 24.41 & 1.45 &
1,307.76 & 7.18 & 7.35 & 0.26 & 1.27 & 1.28 & 49.26 & 45.10 & 49.31 & 0.28 \\
50$\times$200 & TR & 10,000 & 75,774 &
1,604.40 & 533.23 & 128.43 & 0.71 & 241.59 & 0.89 & 37.73 & 1.24 &
1,274.99 & 8.23 & 7.12 & 0.31 & 1.17 & 1.02 & 58.04 & 46.95 & 74.54 & 0.62 \\
50$\times$200 & C & 10,000 & 75,774 &
1,771.62 & 391.80 & 145.47 & 1.43 & 150.05 & 0.64 & 24.88 & 1.12 &
1,297.14 & 13.33 & 7.50 & 0.13 & 1.38 & 1.21 & 52.33 & 46.35 & 54.52 & 0.41 \\
\midrule

100$\times$500 & BL & 50,000 & 389,474 &
12,107.30 & 731.34 & 746.59 & 3.01 & 707.64 & 1.75 & 125.53 & 3.70 &
36,145.90 & 508.87 & 43.66 & 0.30 & 6.84 & 2.98 & 392.65 & 104.07 & 476.38 & 8.54 \\
100$\times$500 & TR & 50,000 & 389,474 &
7,709.86 & 739.94 & 456.42 & 0.62 & 1,251.35 & 62.95 & 187.43 & 3.18 &
33,500.73 & 357.27 & 46.67 & 0.59 & 8.23 & 3.21 & 465.54 & 100.17 & 962.73 & 55.57 \\
100$\times$500 & C & 50,000 & 389,474 &
9,317.94 & 808.51 & 533.08 & 2.85 & 735.58 & 0.94 & 128.85 & 2.56 &
34,934.39 & 305.14 & 50.08 & 0.42 & 7.05 & 3.44 & 414.50 & 110.73 & 560.41 & 20.98 \\
\midrule

100$\times$1000 & BL & 100,000 & 779,474 &
24,709.16 & 841.54 & 1,518.91 & 2.38 & 1,356.78 & 4.52 & 241.18 & 2.32 &
\multicolumn{2}{c}{N/A} & 107.11 & 0.35 & 16.09 & 2.72 & 1,197.74 & 130.81 & 2,364.13 & 82.27 \\
100$\times$1000 & TR & 100,000 & 779,474 &
15,278.99 & 623.56 & 888.21 & 3.24 & 2,456.90 & 81.72 & 374.99 & 4.31 &
\multicolumn{2}{c}{N/A} & 105.45 & 0.30 & 18.67 & 2.28 & 1,493.55 & 81.62 & 4,397.75 & 61.05 \\
100$\times$1000 & C & 100,000 & 779,474 &
18,540.63 & 191.02 & 1,040.46 & 3.83 & 1,402.35 & 4.01 & 246.00 & 3.04 &
\multicolumn{2}{c}{N/A} & 94.83 & 0.38 & 17.06 & 2.62 & 1,161.87 & 112.41 & 2,745.41 & 11.09 \\
\midrule

100$\times$1500 & BL & 150,000 & 1,169,474 &
36,972.08 & 1,340.97 & 2,296.98 & 8.48 & 2,050.10 & 93.87 & 361.31 & 4.87 &
\multicolumn{2}{c}{N/A} & 150.44 & 0.42 & 23.89 & 2.88 & 2,498.22 & 201.72 & 5,123.66 & 87.66 \\
100$\times$1500 & TR & 150,000 & 1,169,474 &
22,367.87 & 83.31 & 1,318.62 & 2.59 & 3,709.88 & 125.01 & 562.00 & 3.36 &
\multicolumn{2}{c}{N/A} & 171.14 & 0.32 & 29.72 & 2.32 & 4,269.01 & 338.45 & 10,195.76 & 170.68 \\
100$\times$1500 & C & 150,000 & 1,169,474 &
27,219.68 & 238.58 & 1,549.40 & 1.13 & 2,089.71 & 4.88 & 360.45 & 2.82 &
\multicolumn{2}{c}{N/A} & 154.30 & 0.55 & 26.79 & 3.30 & 2,612.86 & 156.94 & 6,164.88 & 52.26 \\
\midrule

100$\times$2000 & BL & 200,000 & 1,559,474 &
49,148.65 & 2,073.36 & 3,073.22 & 11.32 & 2,691.37 & 3.68 & 477.72 & 4.61 &
\multicolumn{2}{c}{N/A} & 197.32 & 0.47 & 32.13 & 2.84 & 5,016.44 & 495.68 & 8,881.13 & 177.23 \\
100$\times$2000 & TR & 200,000 & 1,559,474 &
29,713.07 & 2,141.94 & 1,752.87 & 3.77 & 4,924.44 & 132.85 & 744.97 & 5.65 &
\multicolumn{2}{c}{N/A} & 223.76 & 0.40 & 40.96 & 2.98 & 7,647.16 & 397.67 & 18,529.00 & 154.66 \\
100$\times$2000 & C & 200,000 & 1,559,474 &
35,953.72 & 925.31 & 2,059.25 & 9.67 & 2,798.23 & 111.36 & 481.67 & 5.33 &
\multicolumn{2}{c}{N/A} & 224.18 & 0.79 & 38.98 & 2.70 & 4,683.30 & 426.53 & 10,864.61 & 57.06 \\

\bottomrule
\end{tabular}
}
\parbox{\linewidth}{\tiny
\color{blue}
\textit{Notes:} BL = bottom-left origin; TR = top-right origin;
C = center origin. N/A indicates that Vanilla Dijkstra was not
evaluated for that instance.
}
\end{table}

\begin{table}[!htp]
\caption{\textbf{Runtime Comparison on Geospatial Graphs
(Mean $\pm$ SD over 10 runs; ms)}}
\label{tab:geospatial}
\centering
\resizebox{\linewidth}{!}{%
\scriptsize\color{blue}
\begin{tabular}{
    l l r r
    r@{$\,\pm\,$}l
    r@{$\,\pm\,$}l
    r@{$\,\pm\,$}l
    r@{$\,\pm\,$}l
    r@{$\,\pm\,$}l
    r@{$\,\pm\,$}l
    r@{$\,\pm\,$}l
    r@{$\,\pm\,$}l
    r@{$\,\pm\,$}l
}
\toprule
\raisebox{0.5\normalbaselineskip}{\textbf{Graph}} &
\raisebox{0.5\normalbaselineskip}{\textbf{Case}} &
\raisebox{0.5\normalbaselineskip}{\textbf{Nodes}} &
\raisebox{0.5\normalbaselineskip}{\textbf{Edges}} &
\multicolumn{2}{c}{\raisebox{0.5\normalbaselineskip}{\textbf{BMSSP-CD}}} &
\multicolumn{2}{c}{\textbf{\shortstack[c]{BMSSP-CD\\(C++)}}} &
\multicolumn{2}{c}{\raisebox{0.5\normalbaselineskip}{\textbf{BMSSP}}} &
\multicolumn{2}{c}{\textbf{\shortstack[c]{BMSSP\\(C++)}}} &
\multicolumn{2}{c}{\textbf{\shortstack[c]{Vanilla\\Dijkstra}}} &
\multicolumn{2}{c}{\textbf{\shortstack[c]{SCGraph\\Dijkstra}}} &
\multicolumn{2}{c}{\textbf{\shortstack[c]{SCGraph Dijkstra\\(C++)}}} &
\multicolumn{2}{c}{\textbf{\shortstack[c]{NetworkX\\Dijkstra}}} &
\multicolumn{2}{c}{\textbf{\shortstack[c]{igraph\\Dijkstra}}} \\
\midrule

Marnet & LA & 11,062 & 34,698 &
810.83 & 328.65 & 49.75 & 0.48 & 212.67 & 0.98 & 29.44 & 0.99 &
1,663.45 & 24.23 & 8.49 & 1.29 & 1.76 & 0.37 & 45.05 & 45.45 & 55.47 & 2.26 \\
Marnet & NY & 11,062 & 34,698 &
804.43 & 281.46 & 48.90 & 0.60 & 219.13 & 3.83 & 28.96 & 0.51 &
1,667.02 & 33.47 & 7.19 & 0.10 & 1.74 & 0.35 & 44.26 & 45.28 & 59.24 & 0.64 \\
Marnet & SEA & 11,062 & 34,698 &
710.50 & 3.65 & 48.85 & 0.21 & 207.56 & 0.76 & 27.83 & 0.32 &
1,636.02 & 7.76 & 7.24 & 0.06 & 1.77 & 0.44 & 49.21 & 53.80 & 69.31 & 0.87 \\
\midrule

North America Rail & LA & 9,929 & 28,915 &
636.63 & 261.44 & 38.28 & 0.33 & 218.36 & 0.98 & 28.87 & 0.33 &
1,289.22 & 4.16 & 5.89 & 0.09 & 1.36 & 0.24 & 39.93 & 45.24 & 52.51 & 0.63 \\
North America Rail & NY & 9,929 & 28,915 &
621.27 & 274.94 & 36.41 & 0.41 & 230.26 & 0.76 & 29.14 & 0.22 &
1,276.90 & 2.52 & 5.76 & 0.03 & 1.35 & 0.26 & 41.65 & 48.25 & 56.16 & 0.22 \\
North America Rail & SEA & 9,929 & 28,915 &
628.45 & 268.66 & 37.29 & 0.36 & 241.48 & 45.82 & 29.10 & 0.18 &
1,279.13 & 1.31 & 5.77 & 0.05 & 1.35 & 0.26 & 39.50 & 43.45 & 48.64 & 0.30 \\
\midrule

Oak Ridge Maritime & LA & 10,661 & 25,036 &
568.26 & 264.02 & 35.59 & 0.39 & 191.42 & 0.38 & 25.17 & 0.42 &
1,464.56 & 15.11 & 5.78 & 0.05 & 1.51 & 0.42 & 36.61 & 51.34 & 31.66 & 0.81 \\
Oak Ridge Maritime & NY & 10,661 & 25,036 &
555.92 & 250.58 & 34.87 & 0.42 & 186.95 & 0.60 & 24.75 & 0.86 &
1,478.70 & 12.10 & 5.57 & 0.03 & 1.44 & 0.38 & 35.29 & 42.94 & 39.53 & 0.32 \\
Oak Ridge Maritime & SEA & 10,661 & 25,036 &
480.16 & 1.38 & 35.41 & 0.41 & 192.69 & 0.46 & 25.44 & 0.53 &
1,457.89 & 7.87 & 5.67 & 0.11 & 1.44 & 0.34 & 39.72 & 48.55 & 52.68 & 0.14 \\
\midrule

US Freeway & LA & 14,591 & 44,232 &
760.26 & 248.77 & 44.73 & 0.34 & 349.40 & 45.25 & 39.35 & 0.39 &
2,819.43 & 30.46 & 7.13 & 0.10 & 1.76 & 0.46 & 78.79 & 46.35 & 186.30 & 2.73 \\
US Freeway & NY & 14,591 & 44,232 &
771.52 & 262.13 & 45.07 & 0.40 & 352.63 & 45.68 & 39.92 & 0.48 &
2,823.42 & 41.62 & 7.19 & 0.20 & 1.67 & 0.41 & 69.31 & 43.45 & 148.48 & 2.04 \\
US Freeway & SEA & 14,591 & 44,232 &
776.31 & 252.09 & 46.01 & 0.25 & 348.28 & 46.70 & 39.78 & 0.45 &
2,814.11 & 18.30 & 7.37 & 0.38 & 1.67 & 0.40 & 88.15 & 44.69 & 218.86 & 5.71 \\
\midrule

World Highways and Marnet & LA & 572,009 & 1,689,541 &
43,003.08 & 1,334.78 & 2,801.33 & 14.15 & 14,639.99 & 178.77 & 2,335.98 & 40.59 &
\multicolumn{2}{c}{N/A} & 476.68 & 1.97 & 128.20 & 1.38 & 8,261.05 & 462.97 & 17,028.90 & 568.46 \\
World Highways and Marnet & NY & 572,009 & 1,689,541 &
43,134.20 & 1,088.52 & 2,830.17 & 3.55 & 15,287.75 & 212.13 & 2,068.67 & 14.77 &
\multicolumn{2}{c}{N/A} & 468.30 & 0.88 & 122.38 & 1.96 & 6,470.21 & 452.73 & 16,459.89 & 177.08 \\
World Highways and Marnet & SEA & 572,009 & 1,689,541 &
43,372.41 & 1,050.39 & 2,878.93 & 8.53 & 14,514.57 & 164.47 & 2,308.73 & 46.24 &
\multicolumn{2}{c}{N/A} & 481.25 & 1.70 & 125.96 & 3.22 & 5,530.60 & 350.42 & 12,222.92 & 172.31 \\

\bottomrule
\end{tabular}
}
\parbox{\linewidth}{\tiny
\color{blue}
\textit{Notes:} LA = Los Angeles origin; NY = New York origin;
SEA = Seattle origin. N/A indicates that Vanilla Dijkstra was not
evaluated for that instance.
}
\end{table}

Three main findings emerge. First, Python BMSSP becomes increasingly competitive with the general-purpose \texttt{NetworkX} and \texttt{igraph} implementations as graph size increases. As shown in Tables~\ref{tab:square_gridgraphs} and \ref{tab:rect_gridgraphs}, the Python implementation of BMSSP outperforms both \texttt{NetworkX} and \texttt{igraph} on the largest square and rectangular gridgraphs. On the largest tested grids, BMSSP is approximately $1.6$--$2.1\times$ faster than \texttt{NetworkX} and $2.4$--$4.2\times$ faster than \texttt{igraph}. This trend is also visible in Panels (d)--(f) of Figure~\ref{fig:algorithm_runtimes}, where the relative performance of BMSSP improves as graph size increases.

Second, implementing BMSSP in C++ substantially reduces its runtime. Across most instances in Tables~\ref{tab:square_gridgraphs}--\ref{tab:geospatial}, BMSSP C++ is approximately $6$--$8\times$ faster than the Python implementation. Panels (g)--(i) of Figure~\ref{fig:algorithm_runtimes} show this comparison directly by expressing all runtimes relative to BMSSP C++. This comparison shows that implementation language and low-level engineering contribute substantially to the runtime gap observed for Python BMSSP.

Third, compilation does not alter the overall comparison with \texttt{SCGraph}. \texttt{SCGraph} Dijkstra C++ achieves the lowest runtime on every tested instance, while its Python implementation is generally the fastest Python baseline, particularly on the larger graphs. BMSSP C++ therefore substantially reduces the implementation overhead observed in Python BMSSP but remains slower than the corresponding \texttt{SCGraph} Dijkstra implementations.

Additional discussion of these performance patterns, including the effect of
the constant-degree transformation, the Python--C++ comparison, and the
remaining performance gap with \texttt{SCGraph}, is provided in
Appendix~\ref{app:runtime_discussion}.
}


{\color{blue}
\subsection{Extrapolation Analysis: BMSSP vs.\ \texttt{SCGraph} Dijkstra}
\label{subsec:break}

To examine whether \texttt{SCGraph} Dijkstra's observed runtime advantage could reverse at larger scales, we fit empirical power-law models \(T(x)=ax^b\) separately for the Python and C++ implementations, where graph size is \(x=n+m\). The analysis uses all 69 benchmark observations, spanning \(1.71\times10^3 \leq x \leq 8.97\times10^6\), with parameters estimated by ordinary least squares on log runtime. We use empirical scaling rather than worst-case complexity because the benchmark graphs are not worst-case instances and implementation constants materially affect observed runtime.

\begin{figure}[!ht]
    \centering
    \caption{Power-law extrapolation of BMSSP and \texttt{SCGraph} Dijkstra
    for the Python (left) and C++ (right) implementations. The intersections
    denote estimated crossovers of \(x^\star\approx5.08\times10^{22}\)
    for Python and \(x^\star\approx4.05\times10^{17}\) for C++.
    Shaded regions show 95\% bootstrap confidence intervals based on
    10,000 paired replications.}
    \includegraphics[width=\linewidth]
    {figures/powerlaw_crossovers_1x2.pdf}
    \label{fig:crossover_uncertainty}
\end{figure}
%\vspace{-20pt}
Figure~\ref{fig:crossover_uncertainty} shows the fitted models. For Python, \(T_{\mathrm{BMSSP}}(x)\approx9.91\times10^{-3}x^{1.0305}\) and \(T_{\mathrm{SCGraph}}(x)\approx6.41\times10^{-4}x^{1.0829}\), with log-scale \(R^2\) values of 0.9961 and 0.9966. The smaller fitted exponent for BMSSP implies a crossover at \(x^\star\approx5.08\times10^{22}\). A future crossover occurs in all 10,000 bootstrap replications at configuration level, with a 95\% confidence interval of \([1.66\times10^{20},\,1.10\times10^{26}]\); even its lower bound lies more than 13 orders of magnitude beyond the largest observed graph.

For C++, \(T_{\mathrm{BMSSP}}(x)\approx5.65\times10^{-4}x^{0.9582}\) and \(T_{\mathrm{SCGraph}}(x)\approx1.08\times10^{-5}x^{1.0557}\), with log-scale \(R^2\) values of 0.9613 and 0.9733. The estimated crossover, \(x^\star\approx4.05\times10^{17}\), is roughly five orders of magnitude closer to the observed range than for Python, but remains far beyond it. A future crossover again occurs in all bootstrap replications, with a 95\% confidence interval of \([2.21\times10^{14},\,1.44\times10^{25}]\).

Thus, both implementations exhibit a smaller fitted scaling exponent for BMSSP, but the estimated break-even point depends strongly on implementation language and remains far outside the observed range. The wide confidence intervals further caution against interpreting these extrapolations as practically attainable thresholds. Appendix~\ref{app:runtime_ratio} examines sensitivity to alternative functional forms and runtime-ratio extrapolations.
}


{\color{blue}
\subsection{Precision Parameter Sensitivity}
\label{subsec:precision_sensitivity}

The BMSSP recursion parameters \(k\), \(t\), and \(l\) are determined by graph size according to the theoretical construction and are therefore not treated as tuning parameters in \texttt{BMSSPy}. The implementation additionally uses the precision parameter \(p\), described in Section~\ref{sec:unique_distances}, which specifies the number of decimal digits retained when constructing deterministic effective distances. We therefore examine whether \(p\) affects the stability of computed shortest-path distances or runtime.

\begin{figure}[!ht]
    \centering
    \caption{Sensitivity of BMSSP to precision \(p\).
    Panels (a)--(e) report absolute percentage differences in selected
    shortest-path distances relative to \(p=8\) for three gridgraphs and
    two geospatial graphs. Panel (f) reports average BMSSP solve time.
    Each precision setting is evaluated over 10 runs.}    
    \includegraphics[width=\linewidth]{figures/precision_benchmark_plots (1).pdf}
    \label{fig:precision_sensitivity}
\end{figure}
%\vspace{-20pt}
We vary \(p\in\{0,2,4,6,8\}\) on three gridgraphs and two geospatial graphs. Figure~\ref{fig:precision_sensitivity} reports absolute percentage differences in selected shortest-path distances relative to \(p=8\), together with average BMSSP solve times over 10 runs. Panels (a)--(e) show that \(p=0\) can introduce substantial error, while \(p=2\) sharply reduces it and \(p=4\), \(p=6\), and \(p=8\) produce indistinguishable distances. Panel (f) shows no runtime increase as precision rises.

These results indicate that precision affects the stability of the computed distances but has little observable effect on runtime over the tested range. Settings of \(p\geq4\) provide stable distances, while very coarse rounding should be avoided. \texttt{BMSSPy} therefore uses \(p=6\) as its default, providing a conservative setting within the empirically stable range. Because appropriate precision may depend on the resolution of the input edge weights, \(p\) remains user-configurable.
}
\section{Conclusion}\label{sec:conc}

The Bounded Multi-Source Shortest Path (BMSSP) algorithm of \citet{duan_breaking_2025} is a major theoretical breakthrough: it breaks the Dijkstra sorting barrier and achieves a deterministic \(O(m\log^{2/3} n)\) time for SSSP on graphs with nonnegative weights. However, {\color{blue}translating this theoretical advance into an executable algorithm requires resolving several implementation details, and existing public implementations do not retain all components required for the stated deterministic complexity bound.}

{\color{blue}
This paper takes a first step toward closing this gap by introducing \texttt{BMSSPy}, to our knowledge, the first publicly released open-source Python package implementing BMSSP at the time of its initial release that retains the algorithmic and data-structural components required for its stated deterministic complexity. Our findings address the three research questions as follows:
}

\begin{itemize}

\item {\color{blue}
\textbf{RQ1.} We can implement BMSSP in a robust and reproducible manner using \texttt{BMSSPy}, which provides an end-to-end implementation that computes shortest-path distances and predecessor paths while retaining BMSSP's recursive structure and required data-structural components. It also handles unreachable nodes and variable node degrees, making the algorithm executable on graphs with nonnegative edge weights.
}

\item {\color{blue}
\textbf{RQ2.} Practical ambiguities can be resolved by making explicit several components left implicit in the theoretical description, including deterministic tie-breaking for non-unique path lengths, predecessor tracking, tight-edge forest and pivot construction, constant-degree conversion, and deterministic set and dictionary operations. The precision analysis further shows that precision \(p\geq4\) produces stable distances on the tested instances, with \(p=6\) used as the package default.
}

\item {\color{blue}
\textbf{RQ3.} BMSSP becomes increasingly competitive with general-purpose Dijkstra implementations as graph size increases: Python BMSSP is faster than \texttt{NetworkX} and \texttt{igraph} on the largest tested gridgraphs, while an equivalent C++ implementation reduces BMSSP runtime by approximately \(6\)--\(8\times\) across most instances. Nevertheless, \texttt{SCGraph} Dijkstra remains faster than BMSSP in both Python and C++ throughout the tested instances. Power-law extrapolations yield potential BMSSP--\texttt{SCGraph} crossovers far beyond the observed graph sizes, with their locations depending strongly on implementation language and model specification; these estimates are therefore better interpreted as scaling diagnostics than as practically attainable break-even points.
}

\end{itemize}

{\color{blue}
These findings highlight the remaining gap between BMSSP's improved worst-case complexity and its current implementation-level performance. \texttt{BMSSPy} therefore currently serves primarily as a reproducible research and reference implementation rather than a drop-in replacement for modern shortest-path solvers/software. Narrowing this gap may involve alternative data structures, bounded-degree handling, work-budget modifications, refined pivot and window policies, or hybrid shortest-path strategies. \texttt{BMSSPy} provides a platform for evaluating whether such algorithm-engineering improvements can reduce this gap.
}

% \section{Conclusion}\label{sec:conc}
% The Bounded Multi-Source Shortest Path (BMSSP) algorithm of \citet{duan_breaking_2025} is a major theoretical breakthrough: it breaks the Dijkstra sorting barrier and achieves a deterministic $O(m\log^{2/3} n)$ time for SSSP on graphs with nonnegative weights. However, until now, this advance has remained largely theoretical, with no public implementation and several key steps left ambiguous in practical settings.


% This paper takes a first step toward closing that gap by introducing \texttt{BMSSPy}, the first functional, open-source Python implementation of BMSSP. Our \text{robust, reproducible implementation} does more than translate pseudocode into code: it addresses potential bugs, relaxes unrealistic assumptions, and makes implicit components, such as tight-edge forest construction, pivot selection, the underlying data structures, and the set operations, fully explicit. We add features such as stable predecessor tracking, handling unreachable nodes, and irregular degree distributions, turning BMSSP into an end-to-end algorithm that returns both distances and routes on real networks.

% % We complement this with \text{comprehensive empirical evaluation} on synthetic gridgraphs and real geospatial networks, comparing BMSSP to Dijkstra implementations in \texttt{NetworkX}, \texttt{igraph}, and \texttt{SCGraph}. BMSSP clearly outperforms a vanilla Python Dijkstra and becomes competitive with general-purpose libraries on large or elongated grids, but SCGraph’s Dijkstra remains substantially faster across all realistic instances. Our break-even analysis shows that the asymptotic advantage of BMSSP would materialize only at astronomically large graph sizes, consistent with the overhead concerns already noted by \citet{duan_breaking_2025}.

% We complement this with comprehensive empirical evaluation on synthetic gridgraphs and real geospatial networks, comparing BMSSP to Dijkstra implementations in \texttt{NetworkX}, \texttt{igraph}, and \texttt{SCGraph}. BMSSP clearly outperforms a vanilla Python Dijkstra and becomes competitive with general-purpose libraries on large or elongated grids, but \texttt{SCGraph}’s Dijkstra remains substantially faster across all realistic instances. Our break-even analysis shows that the asymptotic advantage of BMSSP would materialize only at astronomically large graph sizes, consistent with the overhead concerns already noted by \citet{duan_breaking_2025}.

% While our current implementation prioritizes algorithmic clarity and theoretical fidelity, this choice naturally incurs runtime and memory overheads compared to libraries like \texttt{SCGraph} and \texttt{igraph}. Consequently, \texttt{BMSSPy} currently serves best as a research instrument and a reference ``blueprint" rather than a drop-in replacement for shortest path implementations. However, we view this as a necessary and exciting milestone: by successfully verifying the complex logic and data structures in Python, we lay a solid foundation for future high-performance ports in other languages. We are optimistic that future research and implementations will dramatically reduce memory footprints and constant-factor delays, allowing the algorithm’s asymptotic power to shine in practice.

% Narrowing this gap will require reducing constant factors. This may be achieved through changing the constant degree requirements of the graph, work budget modifications, faster data structures, refined pivot and window policies, or hybrid strategies that combine BMSSP’s bounded-range logic with other heuristics. We hope that \texttt{BMSSPy} will serve as a platform for this next step. The package is fully documented, reproducible, and designed to support future research and implementations.%including time-windowed paths \citep{desaulniers2000shortest}, multi-pair queries \citep{wang2005multiple}, dynamic arc weights \citep{ahuja2003dynamic,buriol2008speeding}.  fixed-charge and resource-constrained variants \citep{engineer2012fixed}, stochastic or risk-averse formulations \citep{bertsekas1991analysis, chicoisne2018risk, bertazzi2024dynamic}, and online environments \citep{lagos2025online}. 
% By providing the first complete, open, and practical implementation of BMSSP, \texttt{BMSSPy} offers a foundation for future algorithmic improvements and makes the BMSSP framework accessible to researchers and practitioners in operations research, transportation, and network analytics.
% These results underscore the gap that often exists between theoretical complexity and system-level performance. Narrowing this gap for BMSSP will require reducing constant factors through simpler data structures, improved cache behavior, refined pivot and window policies, or hybrid strategies that combine BMSSP’s bounded-range logic with practical Dijkstra heuristics. We hope that \texttt{BMSSPy} will serve as a platform for this next step. The package is fully documented, reproducible, and designed to support extensions to both classical and emerging settings, including time-windowed paths \citep{desaulniers2000shortest}, multi-pair queries \citep{wang2005multiple}, dynamic arc weights \citep{ahuja2003dynamic,buriol2008speeding}, fixed-charge and resource-constrained variants \citep{engineer2012fixed}, stochastic or risk-averse formulations \citep{bertsekas1991analysis, chicoisne2018risk, bertazzi2024dynamic}, and online environments \citep{lagos2025online}. By providing the first complete, open, and practical implementation of BMSSP, \texttt{BMSSPy} offers a foundation for future algorithmic improvements and makes the BMSSP framework accessible to researchers and practitioners in operations research, transportation, and network analytics.





% Acknowledgments here
% \ACKNOWLEDGMENT{%
% % Enter the text of acknowledgments here
% }% Leave this (end of acknowledgment)

% --- Appendices (optional) ---
% Appendix here
% Options are (1) APPENDIX (with or without general title) or 
%             (2) APPENDICES (if it has more than one unrelated sections)
% Outcomment the appropriate case if necessary
%
% \begin{APPENDIX}{<Title of the Appendix>}
% \end{APPENDIX}
%
%   or 
%
% \begin{APPENDICES}
% \section{<Title of Section A>}
% \section{<Title of Section B>}
% etc
% \end{APPENDICES}


% References here (outcomment the appropriate case) 
 
% CASE 1: BiBTeX used to constantly update the references 
%   (while the paper is being written).
\clearpage
\bibliographystyle{informs2014} % outcomment this and next line in Case 1
\bibliography{refs/references} % if more than one, comma separated

% CASE 2: BiBTeX used to generate mypaper.bbl (to be further fine tuned)
%\input{mypaper.bbl} % outcomment this line in Case 2

\clearpage % optional, if you want the notation table to start on a fresh page
%\section*{Appendix}
\begin{APPENDICES}
\OneAndAHalfSpacedXI
\textbf{E-Companion for BMSSPy: A Python Package and Empirical Comparison
of Bounded Multi-Source Shortest Path Algorithm}


\section{Installation and Usage Example}\label{app:install_example}

\subsection{Installation}\label{subsec:install}

\texttt{BMSSPy} is available on PyPI and requires Python~3.11 or later:
\begin{lstlisting}[language=bash]
pip install bmsspy
\end{lstlisting}

\subsection{Usage Example}\label{subsec:use}

The following example creates a small directed graph and demonstrates two use cases: computing distances and predecessors from a single source, and extracting a shortest path to a destination. {\color{blue}For a graph with \(n\) nodes, BMSSPy identifies nodes by integer indices \(0,\ldots,n-1\).} The main inputs are:
\begin{itemize}
  \item \texttt{graph}: a list of dictionaries, where \texttt{graph[u]} maps each neighbor \texttt{v} of node \texttt{u} to its nonnegative edge weight \texttt{w}. Undirected graphs can be represented by including each edge in both directions.
  \item \texttt{origin\_id}: an integer node index or a set of indices. If a set is provided, all listed nodes are treated as sources with distance zero.
  \item \texttt{destination\_id}: an integer node index or \texttt{None}. If \texttt{None}, the solver returns distances and predecessors for all reachable nodes.
\end{itemize}

The solver returns a dictionary containing \texttt{origin\_id}, \texttt{destination\_id}, and:
\begin{itemize}
  \item \texttt{distance\_matrix}: shortest-path distances from the origin or origin set; unreachable nodes have distance \texttt{inf}.
  \item \texttt{predecessor}: predecessor indices used for path reconstruction; unreachable nodes have predecessor \texttt{-1}.
  \item \texttt{path}: the sequence of node indices on the shortest path when a destination is specified; otherwise \texttt{None}.
  \item \texttt{length}: the corresponding shortest-path length when a destination is specified; otherwise \texttt{None}.
\end{itemize}

\paragraph{Example.}
The following code builds a small directed graph, visualized in Figure~\ref{fig:graph}, computes shortest-path distances and predecessors from node~0, and extracts the shortest path to node~4. Multi-source queries are specified by providing a set of origin indices. Section~\ref{sec:func} describes the underlying implementation.


\clearpage
\begin{figure}[!ht]
\begin{center}
\caption{Visualization of the 5-node Example Graph}\label{fig:graph}
\includegraphics[width=0.3\linewidth]{figures/graph.pdf}%
\end{center}
\end{figure}

\begin{lstlisting}[language=python]
from bmsspy import Bmssp

# Graph with 5 nodes: 0..4
# Adjacency-list representation with nonnegative weights
graph = [
    {1: 1, 2: 1},   # 0 -> 1 (1), 0 -> 2 (1)
    {2: 1, 3: 3},   # 1 -> 2 (1), 1 -> 3 (3)
    {3: 1, 4: 2},   # 2 -> 3 (1), 2 -> 4 (2)
    {4: 2},         # 3 -> 4 (2)
    {}              # 4 has no outgoing edges
]
bmssp_graph = Bmssp(graph) # Initialize the graph as a Bmssp graph

# Distances and predecessors from origin 0
res_0 = bmssp_graph.solve(origin_id=0) 
print(res_0) #=>
# {
#     'origin_id': 0,
#     'destination_id': None,
#     'predecessor': [-1, 0, 0, 2, 2],
#     'distance_matrix': [0, 1, 1, 2, 3],
#     'path': None,
#     'length': None
# }

# Shortest path from 0 to 4
res_0_4 = bmssp_graph.solve(origin_id=0, destination_id=4) 
print(res_0_4) #=>
# {
#     'origin_id': 0,
#     'destination_id': 4,
#     'predecessor': [-1, 0, 0, 2, 2],
#     'distance_matrix': [0, 1, 1, 2, 3],
#     'path': [0, 2, 4],
#     'length': 3
# }
\end{lstlisting}


%\clearpage

\section{Deterministic Set and Dictionary Operations Details}
\label{app:set_dictionary_details}

This appendix describes the implementation underlying the deterministic set and dictionary operations summarized in Section~\ref{subsec:set_functions}. The construction exploits the fact that every node is represented by an integer index \(0,\ldots,n-1\), allowing membership and value information to be accessed directly by index rather than through hashing.

\subsection{Set Representation}

For each logical set \(S_a\), we maintain three objects:
\begin{itemize}
    \item a list \(S_a\) containing the elements currently stored in the set;
    \item a membership array \(\mathit{memb}_a[0..n-1]\); and
    \item a generation counter \(\mathit{scnt}_a\).
\end{itemize}

The membership array is initialized once to \(-1\). Whenever a new logical instance of \(S_a\) is created, the list is reset to empty and \(\mathit{scnt}_a\) is incremented. An element \(x\) belongs to the current set exactly when
\(
\mathit{memb}_a[x]=\mathit{scnt}_a.
\)

To insert \(x\), we first test this equality. If it is false, \(x\) is appended to \(S_a\) and
\(
\mathit{memb}_a[x]\leftarrow\mathit{scnt}_a.
\)
Because these operations use direct array indexing, membership testing and insertion both require deterministic \(O(1)\) time. Resetting the set does not require clearing the membership array; incrementing the generation counter and resetting \(S_a\) as an empty list instead invalidates all entries from the previous generation. 

\subsection{Set Extension}

To extend \(S_a\) by another set \(S_b\), we iterate through the list representation of \(S_b\). For each \(x\in S_b\), \(x\) is appended to \(S_a\) only if
\(
\mathit{memb}_a[x]\neq\mathit{scnt}_a.
\)
Extending \(S_a\) by \(S_b\) therefore requires \(O(|S_b|)\) time while preserving deterministic \(O(1)\) membership tests. Explicit set removal is not required by the BMSSP implementation.

\subsection{Dictionary Extension}

The same construction supports dictionary-style operations by introducing a direct-indexed value array
\(
V_a[0..n-1].
\)
Assigning value \(y\) to key \(x\) sets
\(
V_a[x]\leftarrow y
\)
and marks \(x\) as active through the membership array. Lookup first checks whether
\(
\mathit{memb}_a[x]=\mathit{scnt}_a;
\)
if so, the corresponding value is returned from \(V_a[x]\). A key can be invalidated by changing its membership marker without clearing the value array.

Thus, lookup, assignment, and invalidation require deterministic \(O(1)\) time. When iteration over active keys is required, the accompanying list \(S_a\) is retained; otherwise, the list can be omitted.

\subsection{State Across Recursion Levels}

Each level of Algorithm~\ref{alg:bmssp_recursive_bmssp} maintains independent set and dictionary state. Separate membership arrays and generation counters are therefore allocated for each recursion level. From Algorithm~\ref{alg:bmssp_solve},
\(
l=\left\lceil\frac{\log_2 n}{t}\right\rceil,
\qquad
t=\left\lfloor(\log_2 n)^{2/3}\right\rfloor,
\)
and hence
\(
l=\Theta((\log n)^{1/3}).
\)
Maintaining \(O(n)\)-length indexed state at each level therefore requires
\(
O\!\left(n(\log n)^{1/3}\right)
\)
one-time initialization and storage.

This construction replaces expected-time hash-table operations with deterministic direct-indexed operations while preserving efficient iteration over active elements.



\section{Data Structure}\label{app:data_structure}

The original data structure contained several ambiguities and open questions that required clarification during implementation. To support the required operations without changing their asymptotic runtime, we introduced several major and minor adjustments that enable the implementation, although some may be unnecessary under alternative designs.

The first major addition is closely related to the mechanisms described in Section~\ref{subsec:set_functions}. We create $l$ recursion level dictionaries that can perform lookup, assignment, and key invalidation. We need $l$ dictionaries since each recursion level has its own state that is independent of the other recursion level states above it. These are used to access the linked list node as well as the appropriate data buffer in $O(1)$ time.


The second major addition is that all linked lists are doubly linked. This includes both the nodes within each block and the blocks within each chain.

We introduce several minor modifications or extensions to support full functionality:
\begin{itemize}
    \item The linked-list blocks referenced by the red–black tree were generalized to chains of linked-list blocks. Although not required when using unique path lengths, this generalization is necessary for supporting non-unique values in future variants. It prevents any block from exceeding size $M$ when more than $M$ keys share the same value. Only the ``head'' block referenced by the red–black tree may contain values less than the lower bound, while all subsequent blocks in the chain contain only values equal to the lower bound. This ensures that pull operations remain $O(M)$ and that insert operations do not need to consider later blocks.
    \item A minimum block size of 2 is enforced, though the return size for the ``pull'' operation may be 1.
    \item Each block’s linked list has an associated container object that tracks its upper bound and current size.
    \item The block containing the global upper bound is never removed from the red–black tree, even if it is empty.
\end{itemize}


\section{Algorithm~\ref{alg:utils_is_pivot}: Is Pivot}
\label{app:is_pivot}
This helper function determines whether a given node qualifies as a pivot by checking the size of its reachable subtree within a directed forest~$F$. 
Using a simple depth-first traversal, it counts visited nodes until reaching the threshold~$k$, returning \textbf{True} if the limit is met and \textbf{False} otherwise. 
This ensures efficient identification of influential frontier roots during pivot selection.
\input{formulas/utils/is_pivot}

\section{Algorithm~\ref{alg:utils_reconstruct_path}: Reconstruct Path}
\label{app:path}

This utility function reconstructs a shortest path from the destination
node~$dest$ to its corresponding origin using the predecessor array
$\textit{pred}$. It follows predecessor links backward until a node with
no predecessor is reached, reverses the resulting sequence, and returns
the ordered path.

\input{formulas/utils/reconstruct_path}
\section{Algorithm~\ref{alg:utils_convert_to_constant_degree}: Convert to Constant Degree}
\label{app:to_out_degree}

This utility function transforms an arbitrary directed graph into an equivalent constant-degree representation in which every node has at most 2 incoming and 2 outgoing edges.
Whenever a node exceeds this limit, it is expanded into a sequence of auxiliary nodes, each responsible for at most one of the original incoming and outgoing edges.
These auxiliary nodes are then linked in a zero-weight cycle to preserve reachability and allow traversal through all incoming and outgoing edges of the original node.
The function returns the augmented graph, a mapping that records the correspondence between new and original nodes, and the size of the original graph.
This conversion enables algorithms that require constant degree without altering the graph’s essential connectivity or path-weight structure.
\input{formulas/utils/convert_to_constant_degree}

\clearpage
\section{Algorithm~\ref{alg:utils_convert_from_constant_degree}: Convert From Constant Degree}
\label{app:from_out_degree}

This utility function reverses the transformation applied during constant-degree conversion, restoring shortest-path data from the augmented graph back to the original graph.
Using the index map produced in the forward conversion, each predecessor entry in the constant-degree graph is traced backward until a valid predecessor from the original graph is identified.
Auxiliary nodes introduced solely to enforce constant degree are skipped transparently, ensuring that no artificial structure influences the recovered paths.
Distances are truncated to the original graph size, and a new predecessor array is constructed in which every predecessor refers only to an original node of the initial graph.
The function returns the reconstructed distance and predecessor matrices, thereby reestablishing shortest-path information in the original graph topology.
\input{formulas/utils/convert_from_constant_degree}

{\color{blue}
\section{Additional Discussion and Details of Runtime Results}
\label{app:runtime_discussion}

This appendix provides additional interpretation of the computational results in Section~\ref{subsec:compar} and Tables~\ref{tab:square_gridgraphs}--\ref{tab:geospatial}. Figure~\ref{fig:algorithm_runtimes_app} reproduces the absolute runtimes from Figures~\ref{fig:algorithm_runtimes}(a)--(c) on the original scale with standard-deviation bands.

\begin{figure}[!ht]
\begin{center}
\caption{Runtime comparison across square gridgraphs, rectangular gridgraphs, and geospatial graphs. Panels (a)--(c) report mean absolute runtimes with standard deviation error bands on the original scale.}
\label{fig:algorithm_runtimes_app}
\includegraphics[width=\linewidth]{figures/algorithm_runtimes_1x3.pdf}
\end{center}
\end{figure}

\paragraph{General-purpose baselines.} BMSSP becomes increasingly competitive with \texttt{NetworkX} and \texttt{igraph} as the gridgraphs grow and outperforms both on the largest tested grids. The pattern is less pronounced on the smaller and structurally heterogeneous geospatial networks, indicating that relative performance depends on both graph size and structure. The slower \texttt{NetworkX} results should also be interpreted in light of its general-purpose design: Python objects, dictionary-based graph storage, flexible weight evaluation, and additional function calls introduce overhead that is largely absent from the more specialized \texttt{BMSSPy} and \texttt{SCGraph} representations. We therefore treat \texttt{NetworkX} as a widely used general-purpose baseline rather than the fastest attainable Python implementation of Dijkstra. Because \texttt{igraph} is C-backed and its internal costs are not instrumented here, we do not attribute its runtime to specific implementation mechanisms.

\paragraph{Python, C++, and \texttt{SCGraph}.} Implementing BMSSP in C++ reduces runtime by approximately \(6\)--\(8\times\) across most instances, showing that language and low-level implementation overhead are substantial. Nevertheless, both Python and C++ \texttt{SCGraph} Dijkstra remain faster than the corresponding BMSSP implementations. On the largest \(1000\times1000\) square grids, BMSSP remains approximately \(13.7\)--\(14.5\times\) slower in Python and \(11.6\)--\(15.8\times\) slower in C++. The persistence of this gap in C++ indicates that it is not solely a Python effect. BMSSP incurs additional work from recursive bounded relaxation, pivot identification, frontier construction, and its supporting data structures, whereas Dijkstra uses a more direct relaxation procedure with efficient priority-queue operations. Some low-level engineering improvements may transfer to BMSSP, but the present experiments do not isolate their individual effects.

\paragraph{Constant-degree transformation.} BMSSP-CD is slower than BMSSP throughout the tested range because the transformation increases the number of nodes and edges. This is especially costly on the gridgraphs, whose degrees are already small and regular. The result does not conflict with the theoretical role of the transformation; it shows only that preserving the constant-degree assumption introduces empirical overhead on these instances.

\paragraph{Overall interpretation.} The results therefore separate three effects: BMSSP can outperform widely used general-purpose implementations on sufficiently large tested grids; compilation substantially reduces BMSSP's runtime; and \texttt{SCGraph} Dijkstra remains faster across the observed benchmark range. The empirical contribution is consequently not runtime dominance, but a reproducible implementation through which the practical behavior and implementation costs of BMSSP can be evaluated directly.


%\clearpage
\section{Extrapolation Scaling Sensitivity and Runtime-Ratio Analysis}
\label{app:runtime_ratio}

The break-even estimates in Section~\ref{subsec:break} rely on a power-law specification and extrapolate far beyond the observed benchmark range. We therefore conduct two complementary robustness analyses. First, we test whether the estimated crossover changes under alternative empirical scaling functions. Second, we examine predicted BMSSP-to-\texttt{SCGraph} runtime ratios over a more moderate extrapolation range. Together, these analyses assess whether the main conclusion depends on the implementation language or the assumed functional form.

\paragraph{Sensitivity to functional form.}
We compare four specifications: \(T=ax^b\), \(T=a(x\log x)^b\), \(T=ax^b(\log x)^c\), and \(\log T=\alpha+b\log x+c(\log x)^2\). Table~\ref{tab:crossover_sensitivity_combined} reports the corresponding fit statistics and implied crossovers for Python and C++. The Python results are particularly sensitive to functional form. The two simpler specifications imply extremely distant crossovers of \(5.08\times10^{22}\) and \(3.03\times10^{23}\), whereas the two more flexible specifications do not predict a BMSSP-faster crossover for \(x\leq10^{100}\). Moreover, the richer specifications provide the strongest fit for Python BMSSP according to both AICc and LOOCV. Thus, for Python, even the existence of a far-future crossover depends on the assumed scaling function.

The C++ results are more consistent in direction but remain uncertain in magnitude. All four specifications predict an eventual crossover, although its location ranges from approximately \(1.37\times10^{14}\) under the log-quadratic model to \(1.36\times10^{18}\) under \(T=a(x\log x)^b\). The two simpler specifications receive the strongest AICc support for both C++ algorithms, while differences in \(R^2\) are small across models. Taken together, these results show that the estimated break-even point depends on both implementation language and functional form. The power-law crossover should therefore be interpreted as an empirical scaling diagnostic rather than as a unique or implementation-independent threshold.

\begin{table}[!ht]
\color{blue}
\centering
\caption{Sensitivity of Python and C++ scaling and crossover estimates to functional form.}
\label{tab:crossover_sensitivity_combined}
\resizebox{\linewidth}{!}{
\begin{tabular}{llccccccc}
\toprule
Implementation & Model
& \(R^2_{\mathrm{B}}\)
& \(R^2_{\mathrm{S}}\)
& \(\Delta\mathrm{AICc}_{\mathrm{B}}\)
& \(\Delta\mathrm{AICc}_{\mathrm{S}}\)
& \(\mathrm{LOOCV}_{\mathrm{B}}\)
& \(\mathrm{LOOCV}_{\mathrm{S}}\)
& Crossover \(x^\star\) \\
\midrule

\multirow{4}{*}{Python}
& \(T=ax^b\)
& 0.9961 & 0.9966
& 19.45 & 0.00
& 0.1550 & 0.1514
& \(5.08\times10^{22}\) \\

& \(T=a(x\log x)^b\)
& 0.9955 & 0.9964
& 29.51 & 4.35
& 0.1672 & 0.1566
& \(3.03\times10^{23}\) \\

& \(T=ax^b(\log x)^c\)
& 0.9972 & 0.9967
& 0.00 & 0.47
& 0.1325 & 0.1507
& -- \\

& \(\log T=\alpha+b\log x+c(\log x)^2\)
& 0.9972 & 0.9967
& 0.27 & 0.55
& 0.1327 & 0.1506
& -- \\

\midrule

\multirow{4}{*}{C++}
& \(T=ax^b\)
& 0.9613 & 0.9733
& 0.12 & 0.00
& 0.4628 & 0.4229
& \(4.05\times10^{17}\) \\

& \(T=a(x\log x)^b\)
& 0.9614 & 0.9733
& 0.00 & 0.04
& 0.4625 & 0.4231
& \(1.36\times10^{18}\) \\

& \(T=ax^b(\log x)^c\)
& 0.9614 & 0.9733
& 2.19 & 2.16
& 0.4674 & 0.4371
& \(3.07\times10^{15}\) \\

& \(\log T=\alpha+b\log x+c(\log x)^2\)
& 0.9614 & 0.9733
& 2.18 & 2.15
& 0.4667 & 0.4342
& \(1.37\times10^{14}\) \\

\bottomrule
\end{tabular}
}
\begin{flushleft}
\footnotesize
B and S denote BMSSP and \texttt{SCGraph}, respectively.
\(\Delta\mathrm{AICc}\) is measured relative to the best specification for each algorithm, so smaller values indicate greater support.
LOOCV reports leave-one-out cross-validation RMSE on the log-runtime scale; smaller values indicate better predictive performance.
A dash indicates that no BMSSP-faster crossover is found for \(x\leq10^{100}\).
\end{flushleft}
\end{table}

\paragraph{Runtime-ratio analysis.}
Because the crossover estimates require extrapolation many orders of magnitude beyond the observed data, we complement them with the runtime ratio
\(
R(x)=\frac{T_{\mathrm{BMSSP}}(x)}{T_{\mathrm{SCGraph}}(x)},
\)
where \(R(x)>1\) indicates that \texttt{SCGraph} is faster. We evaluate this ratio at the largest observed graph, \(x_{\max}=8{,}974{,}074\), and at graph sizes 10, 100, and 1,000 times larger. Table~\ref{tab:runtime_ratio_combined} reports the resulting point estimates and 95\% bootstrap confidence intervals.

For Python, the practical conclusion is stable across specifications. At \(x_{\max}\), the predicted ratio ranges from approximately 6.7 to 7.5. At \(1000x_{\max}\), all point estimates remain well above one, ranging from 4.65 to 15.32, and every 95\% confidence interval remains above break-even. The simpler specifications predict a gradually narrowing gap, whereas the richer models predict that the gap eventually widens, but none approaches a practically relevant crossover over this range.

For C++, all four point estimates also remain above one through \(1000x_{\max}\), but the predicted gap narrows more consistently. At \(x_{\max}\), BMSSP C++ is estimated to require approximately 10.7--11.0 times the runtime of \texttt{SCGraph} C++; at \(1000x_{\max}\), the corresponding point estimates fall to approximately 4.5--5.7. Uncertainty, however, increases rapidly under the more flexible specifications. In particular, the 95\% confidence interval for the log-quadratic model at \(1000x_{\max}\) includes one, so a crossover can no longer be ruled out under that specification.

Overall, the runtime-ratio analysis reinforces the empirical conclusion from Section~\ref{subsec:break}. \texttt{SCGraph} is substantially faster throughout the observed benchmark range in both Python and C++. The Python results preserve this advantage robustly under extrapolation through \(1000x_{\max}\), whereas the C++ results suggest a narrowing performance gap but also substantially greater uncertainty farther beyond the observed data. Thus, neither the power-law crossover estimates nor the alternative scaling models provide evidence that the current BMSSP implementations attain a runtime advantage over \texttt{SCGraph} within the range of graph sizes supported by the experiments.

\begin{table}[!ht]
\centering
\color{blue}
\caption{Predicted BMSSP-to-\texttt{SCGraph} runtime ratios for Python and C++. Entries report point estimates with 95\% bootstrap confidence intervals. Values above one favor \texttt{SCGraph}.}
\label{tab:runtime_ratio_combined}
\resizebox{\linewidth}{!}{
\begin{tabular}{llcccc}
\toprule
Implementation & Model
& \(x_{\max}\)
& \(10x_{\max}\)
& \(100x_{\max}\)
& \(1000x_{\max}\) \\
\midrule

\multirow{4}{*}{Python}
& \(T=ax^b\)
& 6.68 [6.42, 6.90]
& 5.92 [5.58, 6.22]
& 5.25 [4.85, 5.62]
& 4.65 [4.22, 5.08] \\

& \(T=a(x\log x)^b\)
& 6.69 [6.43, 6.90]
& 5.95 [5.61, 6.24]
& 5.29 [4.90, 5.65]
& 4.71 [4.28, 5.12] \\

& \(T=ax^b(\log x)^c\)
& 7.36 [7.22, 7.53]
& 7.70 [7.38, 8.08]
& 8.39 [7.80, 9.10]
& 9.45 [8.48, 10.64] \\

& \(\log T=\alpha+b\log x+c(\log x)^2\)
& 7.50 [7.34, 7.74]
& 8.53 [8.03, 9.18]
& 10.82 [9.60, 12.39]
& 15.32 [12.57, 19.00] \\

\midrule

\multirow{4}{*}{C++}
& \(T=ax^b\)
& 10.95 [9.63, 12.48]
& 8.74 [7.10, 10.82]
& 6.97 [5.23, 9.40]
& 5.57 [3.85, 8.16] \\

& \(T=a(x\log x)^b\)
& 11.00 [9.67, 12.53]
& 8.84 [7.19, 10.91]
& 7.11 [5.34, 9.52]
& 5.73 [3.98, 8.33] \\

& \(T=ax^b(\log x)^c\)
& 10.75 [8.97, 12.52]
& 8.31 [5.03, 12.40]
& 6.37 [2.52, 13.21]
& 4.85 [1.14, 14.95] \\

& \(\log T=\alpha+b\log x+c(\log x)^2\)
& 10.73 [8.78, 12.73]
& 8.20 [4.23, 14.05]
& 6.14 [1.56, 18.57]
& 4.51 [0.44, 29.42] \\

\bottomrule
\end{tabular}
}
\begin{flushleft}
\footnotesize
\(x_{\max}=8{,}974{,}074\), so \(1000x_{\max}\approx8.97\times10^9\).
Confidence intervals are based on 10,000 paired bootstrap replications.
\end{flushleft}
\end{table}
}



\end{APPENDICES}

\end{document}